% ch-verification.tex -- verification: every claim, how strong it is, and how to reproduce it
\chapter{Verification}
\label{ch:verification}

This chapter is the evidence. It records what has been checked in Yuzen, what
the checks measured, and what the numbers came out at. It does not re-derive
anything: the governing equations are in \cref{ch:governing}, the physical
terms in \cref{ch:terms}, the wetting model in \cref{ch:wetting} and the
numerical method in \cref{ch:discretization,ch:timeintegration,ch:convergence}.
What is new here is the evidence and only the evidence.

Two habits run through every section. \textbf{Each claim says how strong it
is}, because ``verified'' covers everything from an algebraic identity to a
number somebody once wrote down and has been comparing against ever since.
\textbf{Each claim says how to reproduce it}, and
\cref{sec:verification:reproduce} collects every command in one place. If a
number in this chapter matters to you, run the command and get it again; if you
cannot, that is a bug in this chapter.

\section{How to read this chapter}
\label{sec:verification:policy}

Claims about a solver come in strengths, and the differences between them are
large. In descending order:

\begin{description}
  \item[Agreement with an analytic solution.] The strongest. There is an
    answer that exists independently of the code, and the code reproduces it.
  \item[Convergence order.] Nearly as strong, and in some ways more
    informative: it says not only that the error is small but that it is the
    error you expect from the discretization you wrote. A wrong coefficient or
    a wrong sign survives a single-mesh comparison far more easily than it
    survives a refinement study.
  \item[Invariants and identities.] Properties that must hold exactly ---
    a sum that must be a plain sum, a uniform state that must stay uniform, a
    reduction that must give the same bits whatever the thread count. These
    are cheap to test and they fail loudly.
  \item[Regression fixing.] The weakest. The current behavior is recorded, and
    a change that alters it is caught. This says nothing whatever about the
    current behavior being \emph{right}; it only says it has not changed.
\end{description}

It is worth naming which of Yuzen's claims sit where, because they are not all
in the same place.

\subsection*{The strong claims}

\begin{itemize}
  \item \textbf{The disjoining-pressure coefficient $\Dpn$ reproduces the
    equilibrium contact angle you ask for.} This is checked in closed form ---
    no numerics at all --- over 36 parameter combinations
    (\cref{sec:verification:disjoining}). The calibration is an identity in the
    limit of an infinite precursor, and the code agrees with the identity to
    within $10^{-14}$ relative.
  \item \textbf{The discretization of capillary plus disjoining pressure is
    second-order accurate.} On an analytic equilibrium droplet the residual
    velocity falls to zero at order 1.94 under mesh refinement
    (\cref{sec:verification:disjoining}). If any of the coefficient, the sign
    or the discretization were wrong, that residual would not vanish under
    refinement at all.
  \item \textbf{The reductions and BiCGStab are bit-identical under a change of
    thread count} (\cref{sec:verification:linalg}). Not ``agree to a tolerance'':
    equal, tested with \code{==}.
\end{itemize}

\subsection*{The weak claims}

\begin{itemize}
  \item \textbf{Cox--Voinov} (\cref{sec:verification:spreading}). What is
    checked is the \emph{form} of the relation and the correlation of the fit.
    The logarithmic factor $\ln(L/l)$ is not independently verified, and the
    agreement with a naive estimate of it claims nothing beyond order of
    magnitude.
  \item \textbf{The cost of the coupling} (\cref{sec:verification:coupling}).
    One 1-D, well-mixed case family, measured in iteration counts. It is not a
    statement about two dimensions, about the \val{"CZ"} model, or about
    wall-clock time.
  \item \textbf{Input interpretation, output format and time-step control.}
    These are regression fixes. They have no external ground truth and make no
    claim to be correct, only to be unchanged.
\end{itemize}

What has not been checked at all is a separate chapter,
\cref{ch:notverified}. Read it before you rely on a part of Yuzen this chapter
does not mention.

\section{The test suite}
\label{sec:verification:suite}

The unit tests build and run with three commands.

\begin{shell}
cmake -S . -B build -DCMAKE_BUILD_TYPE=Release
cmake --build build -j
ctest --test-dir build
\end{shell}

\code{ctest} runs the same executable three times, in separate processes: once
at the default thread count, once with \env{OMP\_NUM\_THREADS} set to \val{1},
and once with it set to \val{4}; it also runs the fourteen other
entries described in \cref{ch:installation}. That is not redundancy. \textbf{Invariance
under a change of thread count is itself one of the claims}
(\cref{sec:verification:linalg,ch:reproducibility}), so it has to be exercised
rather than assumed.

Calling the executable directly prints the counts.

\begin{plain}
$ ./build/bin/yuzen_tests
[doctest] test cases:    128 |    128 passed | 0 failed | 0 skipped
[doctest] assertions: 183297 | 183297 passed | 0 failed |
\end{plain}

\Cref{tab:testfiles} breaks that down by file.

\begin{yzntable}{The seventeen test files, their cases and their assertions\label{tab:testfiles}}{@{}L{0.28\textwidth} d d L{0.375\textwidth}@{}}
\textsf{\bfseries File} & \textsf{\bfseries Cases} &
\textsf{\bfseries Assertions} & \textsf{\bfseries What it covers} \\ \midrule
\file{test\_grid.cpp}               & 9 & 6,252 & Grid indices, moves and coordinates; where fields sit \\
\file{test\_linear\_solver.cpp}     & 6 & 1,663 & The analytic solution of the 1-D Poisson problem, second-order convergence, bit reproducibility \\
\file{test\_finite\_differences.cpp}& 11 & 171,651 & Interpolation coefficients and gradients against analytic values; every convection scheme: exactness, bounds, symmetry, frozen weights \\
\file{test\_boundary\_condition.cpp}& 11 & 1,025 & Matrix consistency over every combination of boundary conditions; the ghost values they give, inlets on every edge, mirrors and transposes; the boundary rows satisfied exactly after a solve \\
\file{test\_mixed\_velocity.cpp}    & 3 & 216 & That the mixed velocity is the plain sum of the individual terms \\
\file{test\_disjoining.cpp}         & 4 & 179 & Calibration of the disjoining pressure and the residual velocity on the equilibrium shape \\
\file{test\_restart.cpp}            & 12 & 193 & The correspondence between lines of the ASCII restart file and grid points; the marker search of the binary reader; the restart state file; damaged restart files refused \\
\file{test\_config.cpp}             & 11 & 76 & Interpretation of the input file, in TOML and in the \file{.inp} format, the report of unknown entries, and the strict conversions \\
\file{test\_vector.cpp}             & 5 & 34 & Basic operations and bit reproducibility of the reductions \\
\file{test\_positivity\_mobility.cpp}& 11 & 858 & That the face mobility is the positivity-preserving form of Diez \& Kondic; with Navier slip, against a reference without cancellation, continuous at its series, its derivative; zero against a \val{"Pinning"} ghost \\
\file{test\_crs\_matrix.cpp}        & 4 & 28 & Element access in CRS and \code{MatVec} \\
\file{test\_timesteps.cpp}          & 13 & 169 & Time advance and output intervals; the output schedule after a restart; the clock advances by the step solved \\
\file{test\_newton\_convergence.cpp}& 9 & 23 & How the outer iteration exits, and when a stagnation fails the step \\
\file{test\_source.cpp}            & 2 & 15 & The nozzle arm's angle against the integral of its rotation rate \\
\file{test\_opt.cpp}               & 2 & 29 & The \env{YUZEN\_OPT} and \env{YUZEN\_GRAIN} switches parsed \\
\file{test\_thread\_cap.cpp}        & 5 & 65 & The limit on the thread count \\
\file{test\_jacobian.cpp}          & 10 & 821 & Each term's Jacobian against a finite difference of the term, 1-D and 2-D, with every convection scheme, with and without slip; no flux through a wall; the capillary and gravity $A\cdot h$ face by face; each term's face fluxes against its $A\cdot h$ \\ \midrule
Total                               & 128 & 183,297 & \\
\end{yzntable}

Read the assertion column with care. \textbf{It reflects the number of times a
loop inside a case reaches a check}, not the number of distinct properties
being asserted. \file{test\_finite\_differences.cpp} dominates the total
because it checks every convection scheme at every face of many random
profiles; that makes it 94 \% of the assertions and nothing like 94 \% of the
coverage.

\section{Grid and finite differences}
\label{sec:verification:grid}

\file{test\_grid.cpp} builds a \code{RectGrid<3>} with
$N_x, N_y, N_z = 9, 7, 5$ and checks the following at \textbf{every grid point
of all eight site types} (\cref{tab:site-types} of \cref{ch:discretization}):

\begin{itemize}
  \item The linear index and the grid coordinates are inverses of one another.
  \item Neighbors close periodically: $N_x$ moves of one cell return to the
    starting point.
  \item A one-cell move and a half-cell move both return to the starting point
    when reversed, and a half-cell move lands on the correct destination site
    type as well.
  \item The index moves under a half-cell move exactly as the rule says ---
    a move of $+\tfrac12$ cell leaves the index alone when that axis's stagger
    bit is \code{0} and raises it by one when the bit is \code{1}, and a move of
    $-\tfrac12$ cell does the reverse.
  \item The coordinates follow the staggered layout, and nothing degenerates in
    a direction that holds only one line, such as $N_y = 1$.
  \item Fields are zero-initialized and can be read back by index, and the
    components of a vector field sit in different places from one another.
\end{itemize}

The difference side is pinned by \file{test\_finite\_differences.cpp}. Over the
seven fixed schemes --- \val{Upwind\_1st}, \val{Upwind\_2nd}, \val{Central\_2nd},
\val{Upwind\_3rd}, \val{QUICK}, \val{KK} and \val{Central\_4th} --- it fixes
that the interpolation coefficients sum to one, so that a constant is
preserved; that the gradient of a linear function is reproduced exactly; that
the gradient of a quadratic agrees to second order or better; and that the
\code{Node} variant is the average of the \code{Xedge} values on either side of
it. The two limited schemes, \val{Koren} and \val{VanLeer}, are checked with the others in \cref{sec:verification:schemes}.

\section{Linear algebra and determinism}
\label{sec:verification:linalg}

\subsection{The Poisson problem}

\file{test\_linear\_solver.cpp} drives the \val{"eigen"} backend through the
\code{ILinearSolverBackend} interface on the one-dimensional Poisson problem
$u'' = 2$ on $[-1, 1]$ with $u(\pm1) = 0$, whose solution is
$u(x) = x^{2} - 1$. All three preconditioners --- \val{None}, \val{Jacobi} and
\val{ILU} --- converge, with a residual below $10^{-11}$
and a maximum error against the analytic solution at the interior points below
$10^{-4}$ (at $n = 202$, $\Delta x = 0.01$, the measured value is
$2.5\times10^{-5}$). Refining through $n = 52, 102, 202, 402$ \textbf{more than
halves the maximum error every time}, which is the second-order convergence the
discretization promises.

Two further properties are fixed in the same file. A zero right-hand side
returns the trivial solution --- only that, and nothing about how many
iterations it takes to get there: returning in zero iterations was an
optimization in an in-house backend tried during development, not a property
of the problem, and \val{"eigen"} iterates 520 times to the same trivial answer.
And \code{Solve()} restores both the matrix and the right-hand
side to the state they were in before the call --- necessary because the
diagonal scaling of the preconditioner overwrites both
(\cref{sec:linsolve:contract}).

\subsection{Determinism}

\Cref{tab:determinism} lists the properties the unit tests fix. The point of
the table is the word \emph{exactly}: these are comparisons with \code{==}, not
with doctest's \code{Approx}.

\begin{yzntable}{Bit-level determinism fixed by the unit tests\label{tab:determinism}}{@{}L{0.26\textwidth} L{0.17\textwidth} L{0.44\textwidth}@{}}
\textsf{\bfseries Object} & \textsf{\bfseries Condition} &
\textsf{\bfseries Claim} \\ \midrule
\code{Vector} inner product, \code{Sum}, \code{Norm} & $N = 100003$ &
  \textbf{Exactly equal} at 1, 2, 4 and 8 threads (tested with \code{==}, not with \code{Approx}) \\
The \val{"eigen"} backend with ILU & $n = 402$ &
  Every component of the solution agrees at 1, 2, 4 and 8 threads, and so do \textbf{the iteration count and the residual} \\
The \val{"eigen"} backend with ILU & the same problem twice &
  Iteration count, residual and solution agree exactly \\
\end{yzntable}

The reductions fix the order of the additions with a block decomposition that
does not depend on the thread count: at least 512 elements per block, at most
256 blocks (\cref{sec:repro:reduction}). $N = 100003$ is chosen deliberately
--- \textbf{it is not a multiple of 256, so the array does not divide evenly
into blocks and the test straddles the block boundaries} rather than sitting
neatly inside them.

The same property is then measured on the whole solver, with
\file{tools/compare.py}, which reports the maximum difference divided by the
scale of the column. \Cref{tab:wholesolver} gives the result; every entry is
exact agreement. The one backend this does not hold for is \val{"amgcl"},
measured separately in \cref{sec:verification:amgclthreads}.

\begin{yzntable}{Whole-solver reproducibility, maximum difference over column scale\label{tab:wholesolver}}{@{}L{0.52\textwidth} d@{}}
\textsf{\bfseries Comparison} & \textsf{\bfseries Result} \\ \midrule
\file{leveling\_1d\_reg}, 1 thread against 2 threads (\val{"eigen"}) & 0.000e+00 \\
\file{leveling\_1d\_reg}, run twice on 2 threads                   & 0.000e+00 \\
\file{allphys\_2d\_reg}, 1 thread against 2 threads (\val{"eigen"}) & 0.000e+00 \\
\end{yzntable}

\section{Backend agreement}
\label{sec:verification:backends}

The linear solver backend is an input-file choice, made separately for the
thickness equation and for the concentration equation (\cref{ch:linearsolver};
the keys are given in \cref{sec:verification:reproduce}).
\file{leveling\_1d\_reg} --- $N_x = 300$, 20 output intervals --- was solved
with \val{"eigen"} and with \val{"amgcl"}, on one thread, and the outputs
compared.
\Cref{tab:backendagreement} gives the result.

\begin{yzntable}{The two backends on \file{leveling\_1d\_reg}, maximum difference over column scale\label{tab:backendagreement}}{@{}L{0.33\textwidth} L{0.24\textwidth}@{}}
\textsf{\bfseries What is compared} &
\textsf{\bfseries \val{"eigen"} vs \val{"amgcl"}} \\ \midrule
Physical quantities in \file{field*.dat} & 7.033e-11 (\code{Ux}) \\
\file{History.dat}                       & 9.615e-02 (\code{ResinVolErr}) \\
Time and step in \file{CalcInfo.dat}     & 0.000e+00 \\
\end{yzntable}

Two readings matter more than the numbers themselves.

\textbf{The $7.033\times10^{-11}$ is the linear solver's tolerance, not a
time-integration discrepancy.} \key{Epsilon} is \val{1e-13} in this case, and a
difference of that order in the solution is what a difference of that order in
the convergence criterion produces. It is \emph{not} an $O(\Delta t)$
disagreement between two different time integrations; the two backends are
solving the same discrete system and stopping at slightly different points
inside the same tolerance.

\textbf{\code{VolErr}, and the conservation-error columns generally,
show a larger relative difference because they are themselves near-zero
quantities.} Their absolute values hover around zero and dividing by their
scale magnifies whatever is left. This case has symmetric edges, so nothing
flows out, and \code{VolErr} and \code{ResinVolErr} are rounding, at most
$7\times10^{-15}$, on which the two backends differ by $6\times10^{-16}$; a
relative difference of $0.096$ on that column does not mean the two backends
disagree about the solute.

The same comparison is what ruled out an in-house BiCGStab backend tried
during development. It agreed with \val{"eigen"} just as closely --- $1.220\times10^{-10}$ in
the physical quantities (\code{Ux}) and $2.935\times10^{-6}$ in the error
measures (\code{ResinVolErr}) --- but it was far more expensive in getting
there. Over the same 20 output intervals the total BiCGStab iteration count was
\textbf{3676} for that backend against \textbf{21} for \val{"eigen"}: a factor
of about 175 for the same answer, which, together with the bugs found in it, is
why it was left out of the release rather than offered as an option. Those totals are summed
from the \code{BCG\_iter} column of \file{CalcInfo.dat} by hand;
\file{tools/compare.py} does not total anything, but it does hold that column,
and the three other diagnostic columns, out of the pass/fail decision and
report their difference separately --- quite correctly, since they are a
property of the solver rather than of the answer.

\subsection{\val{"amgcl"} is not thread-invariant}
\label{sec:verification:amgclthreads}

\val{"amgcl"} is the one backend whose result moves when the thread count
moves. Its internal reductions are not order-fixed the way
\cref{sec:repro:reduction} fixes Yuzen's own, so changing the number of threads
can change the low-order bits. The same \file{leveling\_1d\_reg}, run on one
thread and on four, gives a difference of $6.089\times10^{-11}$ in the physical
quantities (worst column \code{Ux}) and $9.615\times10^{-2}$ in the error
measures (worst column \code{ResinVolErr}, which is rounding here: at most
$6\times10^{-15}$, and the two runs differ by $6\times10^{-16}$). On two
threads it gave the same bits as on one, and so did \file{allphys\_2d\_reg}
on one, two and four threads: whether the bits move depends on how the
library splits the work, and nothing guarantees that they do not. For
\val{"eigen"} the identical comparison gives $0.000\mathrm{e}{+}00$, exactly,
at any thread count (\cref{tab:wholesolver}).

\begin{warn}[What the \val{"amgcl"} measurement does and does not cover]
  \code{ctest} passes on an AMGCL build as well, but the agreement above and
  the thread
  differences above come from one regression case run by hand, not from a
  property the test suite re-checks on every commit. The bit-identity that
  \cref{sec:verification:linalg} asserts covers \val{"eigen"}
  only, and \val{"amgcl"} cannot be brought under it --- it is not
  bit-reproducible by construction. See \cref{sec:notverified:numerics}.
\end{warn}

\section{Velocity composition}
\label{sec:verification:velocity}

The mixed velocity is the film's velocity in the output (\code{Ux},
\code{Uy}), and the Courant check is formed from it (the well-mixed
concentration itself is carried by the face volume fluxes, and what crosses
an open edge is counted from them, \cref{sec:disc:wm-conservative}), so if the rule by
which the individual physical terms are combined into it goes wrong, all of
these are wrong. The sign convention that comes from writing the governing equation
in the form $-\nabla\cdot(\dots)$ is already inside each term's
\code{\_Magnitude} (\cref{ch:terms}), which means \textbf{the composition side
must be a plain sum and nothing more}.

\file{test\_mixed\_velocity.cpp} fixes that with powers of two. Six terms are
given distinct values --- capillary pressure 1, Marangoni 2, vertical gravity
4, in-plane gravity 8, centrifugal force 16, disjoining pressure 32 --- and the
mixed velocity must come out at \textbf{63}. Any term dropped, doubled or
counted twice changes the total to a value that identifies it uniquely.

The signs are checked separately, and just as bluntly: vertical gravity alone
set to $+3$ must give a mixed velocity of $+3$, and centrifugal force alone set
to $-5$ must give $-5$. \textbf{Signs pass through unchanged.} In two
dimensions the same holds for the $x$ and the $y$ component alike.

\section{The same problem posed twice}
\label{sec:verification:equivalence}

A test against a known answer needs a known answer; most of the terms have
none (\cref{sec:notverified:physics}). What every term does have is a set of
changes to the input that must not change the physics: exchanging $x$ and
$y$, mirroring $x$, running a film that does not vary in $y$ as 1-D or as
2-D, moving a periodic film round the domain by whole cells, and making the
problem dimensionless with other scales. \file{tools/equivalence.py} runs a
shipped case and a copy of it rewritten in one of those ways, at a fixed time
step, and compares the two results point by point: a term that treats one
direction differently from the other, an edge or a corner handled on one side
only, an index slip in a stencil or in an output writer, or a group of
dimensionless numbers that does not cancel, all show up far above rounding.
\Cref{tab:equivalence} lists the thirty-three checks. \code{ctest} runs them as
\code{equivalence}, in about half a minute.

\begin{yzntable}{The checks of \file{tools/equivalence.py}\label{tab:equivalence}}{@{}L{0.22\textwidth} L{0.52\textwidth} d@{}}
\textsf{\bfseries Check} & \textsf{\bfseries The same problem, posed again as} & \textsf{\bfseries Worst}\\ \midrule
\code{wm2d\_transpose} & $x \leftrightarrow y$ (grid with $\Delta x \ne \Delta y$, shapes, gravity, axis, edges) & $2\times10^{-9}$ \\
\code{wm2d\_mirror\_x} & $x \to -x$ & $1\times10^{-15}$ \\
\code{cz2d\_dry\_transpose} & a drying 2-D \val{"CZ"} film with the Marangoni term, $x \leftrightarrow y$ & $3\times10^{-11}$ \\
\code{cz2d\_dry\_mirror\_x} & the same, $x \to -x$ & $3\times10^{-11}$ \\
\code{wm2d\_spin\_transpose} & the transpose with a strong rotation and UTOPIA for the film height, so that the upwind side of the centrifugal flux is chosen from the cell centers & $8\times10^{-10}$ \\
\code{wm2d\_spin\_mirror\_x} & the mirror, the same & $2\times10^{-9}$ \\
\code{wm2d\_nozzle\_transpose} & the transpose, with a fixed nozzle off the center & $2\times10^{-10}$ \\
\code{wm2d\_nozzle\_thin} & the same on cells 5.6 times longer one way than the other & $2\times10^{-10}$ \\
\code{wm2d\_nozzle\_mirror\_x} & the mirror, with the same nozzle & $2\times10^{-10}$ \\
\code{wm\_1d\_vs\_2d} & 2-D, uniform and periodic in $y$ & 0 \\
\code{wm\_1d\_vs\_2d\_koren} & the same with the Koren scheme for the concentration, whose limiter is settled at the answer of each step and whose explicit half is formed from the state the step starts from & 0 \\
\code{wm\_io\_1d\_vs\_2d} & the same with a \val{"FreshInlet"} and an \val{"Outlet"} in $x$, the film's rim across the inlet & $4\times10^{-10}$ \\
\code{wm\_io\_1d\_vs\_2d\_sym} & and with $y$ symmetric, so the inlet and outlet corners are walls & $9\times10^{-10}$ \\
\code{wm\_io\_1d\_vs\_2d\_ymin} & the inlet run flowing along $+y$ from an inlet on Ymin, $x$ periodic & $9\times10^{-10}$ \\
\code{wm\_io\_1d\_vs\_2d\_ymax\_sym} & along $-y$ from an inlet on Ymax, $x$ symmetric & $3\times10^{-10}$ \\
\code{wm1d\_io\_mirror} & the 1-D inlet run mirrored, inlet at Xmax & $5\times10^{-10}$ \\
\code{wm1d\_shift} & a periodic film moved 13 cells round the domain & $1\times10^{-9}$ \\
\code{wm2d\_shift} & the same in 2-D, 5 cells in $x$ and 7 in $y$ & $5\times10^{-10}$ \\
\code{wm1d\_scales} & \key{TotalHeight} $\times 2.5$ and $L_0 = \SI{0.7}{\milli\metre}$ & $6\times10^{-9}$ \\
\code{cz1d\_hscale} & \val{"CZ"}, \key{TotalHeight} $\times 2.5$ & $1\times10^{-14}$ \\
\code{cz1d\_lscale} & \val{"CZ"}, $L_0 = \SI{0.7}{\milli\metre}$ & $1\times10^{-14}$ \\
\code{wm1d\_nozzle\_mirror} & a nozzle off the center of a flat film, $x \to -x$ & $4\times10^{-12}$ \\
\code{cz1d\_mirror} & \val{"CZ"}, $x \to -x$ & $3\times10^{-14}$ \\
\code{cz1d\_shift} & \val{"CZ"}, a periodic film moved 11 cells & $2\times10^{-13}$ \\
\code{cz2d\_transpose} & \val{"CZ"}, $x \leftrightarrow y$ & $2\times10^{-13}$ \\
\code{cz\_1d\_vs\_2d} & \val{"CZ"}, 2-D uniform in $y$ & $1\times10^{-14}$ \\
\code{cz2d\_deficit\_transpose} & \val{"CZ"}, $x \leftrightarrow y$, restarted with a column next to a wall at $c = -0.05$, which the columns beside it fill & $2\times10^{-15}$ \\
\code{body\_velocity} & flat film: $\rho g h^2/3\mu$ and $\rho\omega^2 r h^2/3\mu$ exactly & $2\times10^{-14}$ \\
\code{cz\_solute} & drying flat \val{"CZ"} film keeps its solute & 0 \\
\code{cz\_solute\_slope} & a drying \val{"CZ"} film over relief keeps its solute & $5\times10^{-15}$ \\
\code{cz\_uniform} & a uniform \val{"CZ"} concentration stays uniform while the film moves & 0 \\
\code{cz\_wm\_advection} & the flow carries the solute as in \val{"WM"} & $2.7\times10^{-3}$ \\
\code{wm\_retry} & a rejected and retried step, against that step taken at its $\Delta t$ & 0 \\
\end{yzntable}

The column \textsf{Worst} is the largest difference over every column of the
field files, relative to how far that field moved over the run; it is the
concentration every time, because the concentration moves least. The film
thickness itself agrees to $2\times10^{-13}$ or better in every well-mixed check and
in the \val{"CZ"} ones without Marangoni flow; with it, the surface-tension
gradient multiplies the rounding of $c$ by $A_s\,\mathrm{Ma}/\Delta x$, the
thickness agrees to $10^{-12}$ and the velocity to $3\times10^{-11}$. The limit the test enforces is $10^{-7}$.

The last six rows are the exceptions to the rule of the section, checks
against a closed form, against the other model or against another run. \code{body\_velocity} reads
the velocity that the in-plane gravity and the centrifugal term write --- the
one that appears in \file{field*.dat} --- on a
flat film in 1-D and 2-D, where it is $\rho g h^2/(3\mu)$ along \key{Acceleration}
and $\rho\omega^2 r h^2/(3\mu)$ away from the axis, and requires it to
$10^{-10}$. \Cref{sec:verification:velocity} checks that the terms are summed
correctly; this checks that each of them has the right size and sign.
\code{cz\_solute} dries a flat \val{"CZ"} film by \SI{0.4}{\percent} and
measures the change of its solute content $h\,\bar c$ as a fraction of the
solvent lost; it is zero to rounding, and the test allows $10^{-10}$. \code{cz\_wm\_advection} switches on a
strong in-plane gravity over a drying film whose evaporation varies along $x$,
with a vertical diffusivity high enough to keep the \val{"CZ"} profile flat,
and compares the change the flow makes to the depth-mean concentration in the
two models. Away from the walls they agree to \SI{0.3}{\percent}; the test
allows \SI{3}{\percent}. This is the check that fixes the direction in which
\val{"CZ"} advects the solute, which none of the symmetries can: a scheme that
carries the solute upstream is as symmetric as one that carries it
downstream. \code{wm\_retry} lets the time-step control reject the first
step of a leveling film 18 times, and compares the step it finally takes
with the same step taken at that $\Delta t$ from the start; the two must agree
to rounding, since both start from the same film. They did not: the retry
built the explicit half of its Crank--Nicolson update from the answer of the
attempt it had just rejected, and was off by $3.5\times10^{-5}$ of the change.
\code{wm1d\_nozzle\_mirror} and \code{wm2d\_nozzle\_mirror\_x} are ordinary
mirrors, but of cases with a nozzle, which no other check has: they caught the footprint being measured
at the end of each sub-division of a cell instead of its center, which put the
supply half a sub-division to one side (\cref{sec:terms:source}).

\file{tools/concz\_vs\_wm.py} poses a film twice in another way, to the two
concentration models (\cref{sec:disc:positivity-concz}), and requires, besides
the film, that the \code{Volume} of their \file{History.dat} agree to
$10^{-6}$. It found that they wrote different quantities for a 1-D film: the
well-mixed model a strip as wide as the length scale $L_0$, \val{"CZ"} one as
wide as \key{Grid.Ymax} $-$ \key{Grid.Ymin} ($8.7\times10^{-12}$ against
$1.0\times10^{-8}$\,\si{\metre\cubed} for the same film). Corrected, both write
the volume of a strip \SI{1}{\metre} wide (\cref{sec:output:history}); they
agree exactly, in 1-D and in 2-D.

\section{Problems with a closed-form answer}
\label{sec:verification:analytic}

\file{tools/analytic.py} sets up problems the lubrication equations answer in
closed form, runs the solver on them and compares. The linear ones start from
a sinusoid of amplitude $10^{-3}$ of a flat \SI{5}{\micro\metre} film on a
periodic domain one wavelength long, $L = \SI{0.25}{\milli\metre}$, and follow
the Fourier coefficient of that wave number over the cell centers: its decay
rate against $\sigma h_0^3 k^4/3\mu$ (capillary), plus $\rho g h_0^3 k^2/3\mu$
with vertical gravity, and its phase against a translation at
$\rho g h_0^2/\mu$ with in-plane gravity. Where the error should fall at a known
rate it is the rate that is checked.

With the disjoining pressure of \cref{sec:terms:disjoining}, a flat film of
thickness $h_0$ is linearly unstable to long waves where $\Pi'(h_0) > 0$: a
sinusoid of wavenumber $k$ grows at
\begin{equation}
  \omega = -\frac{h_0^3}{3\mu}\,k^2\bigl(\sigma k^2 - \Pi'(h_0)\bigr),\qquad
  \Pi'(h) = -\frac{B}{h}\Bigl[n\Bigl(\frac{\hstar}{h}\Bigr)^{n} - m\Bigl(\frac{\hstar}{h}\Bigr)^{m}\Bigr],
  \label{eq:spinodal}
\end{equation}
with $B$ the calibrated coefficient of \cref{sec:terms:disjoining}; for
$(n,m) = (3,2)$ that is $h_0 > 1.5\hstar$. In the dimensionless variables it
reads $\omega = -\tfrac13 H_0^3k^2\bigl(k^2 + \Bo - \Dpn\,\Pihat'(H_0)\bigr)$.
The \code{spinodal} checks follow a sinusoid on a film on either side of
$1.5\hstar$, in the linear range (the unstable one grows 1.85 times).
\Cref{tab:analytic} gives the result;
\code{ctest} runs it as \code{analytic}, in about a minute.

\begin{yzntable}{The checks of \file{tools/analytic.py}\label{tab:analytic}}{@{}L{0.18\textwidth} L{0.58\textwidth} L{0.16\textwidth}@{}}
\textsf{\bfseries Check} & \textsf{\bfseries What is compared} & \textsf{\bfseries Result}\\ \midrule
\code{leveling\_1d} & capillary decay rate, 16 and 32 cells per wavelength & order 1.99 in $\Delta x$ \\
\code{leveling\_2d} & the same for $\cos k_x x\cos k_y y$, 12 and 24 cells & order 1.98 \\
\code{leveling\_dt} & the same rate, $\Delta t$ = 8 and \SI{4}{\milli\second} & order 2.00 in $\Delta t$ \\
\code{leveling\_cz} & \val{"CZ"} against \val{"WM"} on the same problem & $1\times10^{-11}$ \\
\code{leveling\_slip\_1d} & Navier slip, $\lambda = h_0/2$ (\key{Physics.Slip}): decay rate $\sigma(h_0^3/3 + \lambda h_0^2)k^4/\mu$, 16 and 32 cells & order 1.99 in $\Delta x$ ($6.4\times10^{-3}$ at 32) \\
\code{leveling\_slip\_dt} & the same, $\Delta t$ = 3.2 and \SI{1.6}{\milli\second} & order 2.00 in $\Delta t$ \\
\code{energy\_\bk leveling}, \code{\_slip}, \code{\_box}, \code{\_2d} & \key{Output.Energy} on a 30\,\% sinusoid leveling (periodic; with slip; between \val{"Symmetry"} walls with vertical gravity; 2-D): the balance of the run $|\mathcal F(T)-\mathcal F(0)+\int_0^T\mathcal D\,\mathrm dt|$ relative to $\int_0^T\mathcal D\,\mathrm dt$ (the steps' balances summed), and $\mathcal F$ never rising from step to step & order 2.00 in $\Delta t$ ($2\times10^{-4}$ to $5\times10^{-5}$ at \SI{2}{\milli\second} and \SI{1}{\milli\second}, leveling); no rise \\
\code{energy\_\bk dewetting} & the same on a film $5.5\hstar$ thick, spinodally unstable, ruptured onto its precursor at about \SI{10}{\second} and gathered into a drop by \SI{40}{\second}, 128 cells, $\Delta t$ = 0.125 and \SI{0.0625}{\second} & order 1.97 in $\Delta t$; no rise \\
\code{energy\_\bk saturate\_mono} & capillary and disjoining pressure on the film of \code{spinodal\_1d\_unstable}, a 30\,\% sinusoid grown into a steady film over \SI{40}{\second}: $\mathcal F$ never rises and $\mathcal D \ge 0$ (with the mean face mobility the disjoining pressure once took, $\mathcal D$ went negative, $5.9\times10^{-5}$ of its largest value) & 0 \\
\code{eq\_slip} & the equilibrium droplet of \cref{sec:verification:disjoining} ($N_x = 2000$) advanced \SI{0.01}{\second} with slip of \SI{1}{\micro\metre}, five times the precursor film: $\max|\Delta h|/\max h$ & $9.5\times10^{-6}$ ($6.5\times10^{-6}$ without slip) \\
\code{gravity\_1d} & decay rate with vertical gravity, 32 cells & $6\times10^{-3}$ \\
\code{potential\_1d} & \key{WritePotential}: the \code{PHI} column of the initial field (a sinusoid under capillary and disjoining pressure and vertical gravity) against $\sigma\nabla^2 h + \Pi(h) - \rho g h$ & $1.5\times10^{-15}$ \\
\code{leveling\_logout\_\bk times} & \file{leveling\_1d\_logout}: 141 field outputs from $10^{-7}$ to \SI{1}{\second}, 20 a decade, \key{ExactTimes}: $|t - t_i|/t_i$ & $0$ \\
\code{leveling\_logout} & the same: the decay rate at each output, 64 cells & $1.6\times10^{-3}$, as at a fixed step \\
\code{spinodal\_1d\_\bk unstable} & disjoining pressure ($\thetae = \ang{10}$, $(n,m) = (3,2)$) on a film $h_0 = 2\hstar$, a sinusoid \SI{1}{\milli\metre} long: growth rate against \cref{eq:spinodal}, 16 and 32 cells & order 2.00 in $\Delta x$ ($2.3\times10^{-3}$ at 32) \\
\code{spinodal\_1d\_\bk unstable\_dt} & the same, $\Delta t$ = 0.4 and \SI{0.2}{\second} & order 2.00 in $\Delta t$ \\
\code{spinodal\_1d\_stable} & the same on $h_0 = 1.25\hstar$, below $1.5\hstar$: decay rate & order 1.99 ($3.9\times10^{-3}$ at 32) \\
\code{spinodal\_1d\_\bk stable\_dt} & the same, $\Delta t$ = 0.4 and \SI{0.2}{\second} & order 2.01 in $\Delta t$ \\
\code{wave\_1d} & damping added by the upwinded in-plane gravity (UTOPIA), 32 and 64 cells & order 3.25 \\
\code{wave\_1d\_speed} & wave speed $\rho g h_0^2/\mu$ (UTOPIA), 32 cells & $5\times10^{-5}$ \\
\code{wave\_2d\_speed} & the same along \key{Acceleration} in 2-D & $1\times10^{-4}$ \\
\code{wave\_1d\_koren} & the damping with the default Koren scheme, 32 and 64 cells & order 3.41 \\
\code{wave\_1d\_koren\_speed} & its wave speed, 32 cells; order in $\Delta x$ & $6.7\times10^{-4}$; 2.85 \\
\code{wave\_2d\_koren\_speed} & the same in 2-D, 24 cells & $1.1\times10^{-3}$ \\
\code{wave\_cz} & \val{"CZ"} decay rate with in-plane gravity, 32 and 64 cells & order 2.10 in $\Delta x$ \\
\code{wave\_cz\_speed} & \val{"CZ"} wave speed, 32 cells & $6\times10^{-3}$ \\
\code{spin\_1d} & $h_0/\sqrt{1 + 4\rho\omega^2 h_0^2 t/3\mu}$ on a strip of a flat film (an \val{"Outlet"} at each end) & $4\times10^{-7}$ \\
\code{spin\_balance} & a 1-D closed box spinning (\val{"Symmetry"} or \val{"Pinning"} walls; \val{"WM"} at a fixed and an adaptive step, \val{"CZ"} with and without evaporation): $|\code{VolErr}|$, $|\code{ResinVolErr}|$ with the outflow into $y$ counted & $\le 10^{-12}$ ($2\times10^{-15}$) \\
\code{evap\_wm} & drying flat film against the ODE, $\Delta t$ = 10 and \SI{5}{\milli\second} & order 2.00 in $\Delta t$ \\
\code{evap\_cz} & the same with \val{"CZ"} & order 2.16 \\
\code{slab\_cz} & \val{"CZ"} profile against the slab series, $N_z = 20$ & $3\times10^{-4}$ (second order in $N_z$) \\
\code{dry\_cz} & \val{"CZ"} drying at $Eh/D \approx 1$ with $D(c)$, against a reference, $N_z$ = 16 and 32 & order 1.81 in $N_z$ \\
\code{evap\_cap} & evaporation capped at $h/2$ and at the solvent per step: $h_1 = \tfrac23 h_0$, $hc$ kept & $3\times10^{-8}$ \\
\code{evap\_solvent} & the solvent limit $E\Delta t \le h(1-c)$, $c_0 = 0.5$, $E\Delta t/h_0$ = 0.3 and 0.75, 1-D and 2-D & $3\times10^{-8}$ \\
\code{dried\_wm}, \code{dried\_cz} & a film dried out to $c = 1$ at a rate independent of $c$: solute kept, $c \le 1$, also with a lengthened step & $\le 10^{-12}$ \\
\code{conc\_overshoot} & a front from $c = 0.99$ to 0.3 carried by a uniform flow, UTOPIA and Central4 in 1-D, UTOPIA in 2-D, Courant number 0.4: the overshoot past $c = 1$ moved, not cut; $|\code{ResinVolErr}|$, and $c \le 1$ & $\le 10^{-12}$ ($2\times10^{-14}$) \\
\code{conc\_overshoot\_koren} & the same with Koren at a Courant number of 1.5 & $\le 10^{-12}$ ($1\times10^{-15}$) \\
\code{conc\_drying} & a front from $c = 0.9$ to 0.3 on a closed periodic film, carried by the flow and drying by Raoult's law, UTOPIA and Koren on 50 cells, UTOPIA in 2-D on 24: $|\code{ResinVolErr}|$ & $\le 10^{-12}$ ($2\times10^{-15}$) \\
\code{conc\_uniform} & a uniform $c$ in a closed box under a nozzle and a strong in-plane gravity (and, in 1-D, the spin): $\max|c - c_0|$, 1-D and 2-D & $\le 10^{-12}$ ($6\times10^{-16}$) \\
\code{conc\_undershoot}, \code{\_cz} & a front from $c = 0.7$ to 0.005 carried by the flow (UTOPIA, KK, Central4; UTOPIA in 2-D; \val{"CZ"} with UTOPIA): the undershoot below $c = 0$ filled from the cells around, $c \ge 0$, $|\code{ResinVolErr}|$ & $\le 10^{-12}$ ($1\times10^{-14}$) \\
\code{wall\_balance\_cz}, \code{\_solute} & a \val{"CZ"} closed box with every term toward the walls: volume and solute, with and without evaporation & $\le 10^{-12}$ \\
\code{leveling\_adaptive}, \code{\_cz} & leveling with an adaptive step: the decay at the reported times, and the clock against the sum of the steps solved & $5\times10^{-4}$ ($2.9\times10^{-4}$ at \key{Criterion} $=10^{-5}$); $10^{-12}$ \\
\code{inlet\_outlet} & a bump carried out through an \val{"Outlet"}, the film fed through a \val{"FreshInlet"}, 100 cells: $\max|\code{VolErr}|$, $|\code{ResinVolErr}|$ & $\le 10^{-12}$ ($1.3\times10^{-15}$) \\
\code{inlet\_outlet\_left} & the same: the film uniform again after the bump left & $3\times10^{-5}$ \\
\code{solute\_balance} & the same film drying, UTOPIA, Koren, an adaptive step, and without the concentration: $\max|\code{VolErr}|$, $|\code{ResinVolErr}|$ & $\le 10^{-12}$ ($2.2\times10^{-15}$) \\
\code{solute\_balance\_2d} & the same in 2-D, uniform in $y$ (UTOPIA, Koren) and open on all four sides with every kind of open corner & $\le 10^{-12}$ ($1.1\times10^{-14}$) \\
\code{evap\_threshold} & one step of a flat film losing $E\Delta t = 0.3, 0.45, 0.75\,h_0$: the cap from $E\Delta t > h/2$ at either end of the step, in the film and the concentration & $10^{-7}$ \\
\code{evap\_threshold\_2d} & the same in 2-D & $10^{-7}$ \\
\code{negpow\_rim} & flat film, no flow, \val{"negPow"}: thickness lost in one step against the profile, through the blend to 3.9 blend widths past \key{NegPowRmax} & $3\times10^{-15}$ \\
\code{negpow\_rim\_slope} & jump of slope where the blend takes over, relative to the slope & $0.056$ (allowed 0.2) \\
\code{negpow\_rim\_flat} & spread of the rate past \key{NegPowRmax}, relative to its value & $0$ \\
\code{wall\_balance}, \code{wall\_balance\_2d} & a closed box, gravity and rotation toward the walls, four convection schemes, every edge and corner: $|\code{VolErr}|$ & $\le 10^{-12}$ (about $10^{-14}$) \\
\code{wall\_balance\_pinned}, \code{\_pinned\_2d} & the same with \val{"Pinning"} edges & $\le 10^{-12}$ \\
\code{wall\_balance\_slip} & the four boxes above with Navier slip of \SI{3}{\micro\metre} & $\le 10^{-12}$ ($1.1\times10^{-14}$) \\
\code{nozzle\_balance} & a nozzle with an adaptive step, 1-D and 2-D, and \file{nozzle\_1d\_rotating} with its arm turning, the film thickened a hundredfold: \code{VolErr} counts what it supplies & $\le 10^{-12}$ ($1\times10^{-13}$) \\
\code{nozzle\_solute} & the same: \code{ResinVolErr} counts the solute that comes with it & $\le 10^{-12}$ ($1.3\times10^{-13}$) \\
\code{nozzle\_supplied} & the same with the liquid supplied at \key{Conc} (\key{ConcType = "supplied"}): \code{ResinVolErr} counts $\int c_S S$, and \code{VolErr}; on the thick film steps whose Newton correction stagnated are cut & $\le 10^{-12}$ ($2\times10^{-13}$) \\
\code{nozzle\_mixing} & liquid at \key{Conc} on a film that does not flow, over the whole domain and under a small nozzle, 1-D and 2-D: $c$ against \cref{eq:source-mixing}; with \val{"local"}, $c = c_0$ & $\le 10^{-12}$ ($4.2\times10^{-15}$) \\
\code{nozzle\_arm} & where a rotating nozzle puts the liquid: steady turn and stop & $2\times10^{-4}$ of the arm \\
\code{nozzle\_area} & a fixed 2-D nozzle on near-square cells and on cells 11 and 20 times longer one way than the other: what it supplies against \key{InletRate} $\pi R^2 t$ & $3.6\times10^{-3}$ (its area from $20\times20$ points a cell) \\
\end{yzntable}

The errors of the linear checks are relative to the rate or the phase that is
compared; those of the flat films are relative to the change over the run.

One thing the table records is a property of the scheme rather than a
confirmation of it: \textbf{the in-plane gravity of \val{"WM"} is upwinded},
and the upwinding damps a traveling wave: with UTOPIA \SI{3.3}{\percent} of
the capillary rate at 32 cells per wavelength, \SI{0.35}{\percent} at 64,
falling at third order (with KK, three times as much: \SI{11}{\percent} and
\SI{1.3}{\percent}). The default Koren scheme damps it by \SI{4.6}{\percent}
and \SI{0.43}{\percent}, but clips the crests and troughs, so its speed error
is 13 times UTOPIA's. \val{"CZ"} takes the same flux centered, and its error falls at
second order from \SI{0.6}{\percent} at 32 cells.

\code{slab\_cz} dries a flat \val{"CZ"} film slowly enough ($E h/D = 5\times
10^{-5}$) for the profile through the thickness to be that of a slab with a
constant flux $q = cE$ into its top face and none through the bottom,
\begin{equation*}
  c - \bar c = \frac{q h}{D}\left[\frac{3\zeta^2 - 1}{6}
    - \frac{2}{\pi^2}\sum_{n\ge1}\frac{(-1)^n}{n^2}\,
      \mathrm{e}^{-n^2\pi^2 D t/h^2}\cos n\pi\zeta\right],
\end{equation*}
and compares the computed profile with it; the error, relative to a third of
$qh/D$, is $1.3\times10^{-3}$, $3.2\times10^{-4}$ and $8.7\times10^{-5}$ at
$N_z = 10$, 20 and 40, falling at second order until the linearization, which
is of order $Eh/D$, sets a floor.

\code{dry\_cz} is the regime \code{slab\_cz} stays out of: a film drying
fast enough ($Eh/D \approx 1$) for the solute to pile up under the surface,
with a diffusivity that falls by $e^{-3}$ per unit of concentration. There is
no closed form, so \file{tools/analytic.py} solves the same problem itself
--- the solute per cell conserved in $\zeta$, $(hc)_t = \partial_\zeta[h_t\zeta c
+ (D/h)c_\zeta]$, with no flux through either face, on 600 cells, implicit
midpoint in time --- and the \val{"CZ"} profile must approach it at second
order in $N_z$ (1.79, 1.81 and 1.87 from $N_z = 8$ to 64, the test asking
for at least 1.8 from 16 to 32). It is the check on the vertical advection that the $\zeta$
mapping produces, and on its upwinding at the surface, which nothing else
constrains.

\code{evap\_cap} dries a flat film so fast that one step would take more than
half of it. The evaporation is then capped at $h/(2\Delta t)$, at the start of
the step and at its end, so the Crank--Nicolson step gives $h_1 = h_0 -
(h_0 + h_1)/4$, $h_1 = 0.6\,h_0$, and the solute $hc$ must stay what it was.
The check found that it did not: the concentration equation enriched the
solute at the uncapped rate and made \SI{30}{\percent} more of it in the one
step (found; corrected). The well-mixed concentration equation is the balance
of $hc$ itself, which the evaporation does not enter: the solvent the
film-height equation removed is in $h_1$, so $hc$ is kept exactly
(\cref{sec:disc:wm-conservative}); with the solvent limit as
well, $h_1 = \tfrac23 h_0$ for $c_0 = 0.5$. \code{evap\_threshold} checks where the cap starts: one step each
below it, across it within the step, and above it at both ends, in the film
and in both concentration equations. An enrichment $c_t = Ec/h$ at the rate
of the field alone kept the solute below the cap only to the order of
the scheme, $1 + O((E\Delta t/h)^3)$ per step, which is what
\code{evap\_wm} measures (found; corrected the same way). \code{negpow\_rim} compares the rate of the
\val{"negPow"} profile with the definition across the rim. It found that the
cosine blending the profile into a constant was never stopped at
\key{NegPowRmax}: beyond it the rate fell again and oscillated, down to 0.75 of
its rim value two blend widths out (at the corners of a square domain whose
half-width is \key{NegPowRmax}, for example). Corrected, it is constant
beyond \key{NegPowRmax}.

\code{inlet\_outlet} feeds a film on an incline through a \val{"FreshInlet"}
and drains it through an \val{"Outlet"}, and carries a bump of
\SI{1.5}{\micro\metre} out of the domain. The volume in the domain plus what
has flowed out must stay what it was, and so must the solute
(\file{History.dat} reports the balances as \code{VolErr} and
\code{ResinVolErr}), and once the bump has gone the film must be uniform
again. \code{wm\_io\_1d\_vs\_2d} of \cref{sec:verification:equivalence} adds
the corners these edges make.

\code{solute\_balance} lets the same film dry while it flows, so its
concentration rises downstream and solute diffuses through both edges: the
inlet holds $c$ at its face, the outlet extrapolates it. The film-height and
the concentration equations are finite-volume balances
(\cref{sec:disc:wm-conservative}), so what is in the domain changes by what
crosses the edge faces and the sources alone, and \file{History.dat} counts
what crosses an open edge from the very face fluxes each committed step put
through it, at the step's weighting in time (\cref{sec:output:history}): the
largest $|\code{VolErr}|$ and $|\code{ResinVolErr}|$ over the run must stay
at rounding, $10^{-12}$, with UTOPIA and with Koren for the concentration
(the step completed with the limiter settled), with the time step cut back
and lengthened by the control (only the committed steps count), and without
the concentration (the volume alone); \code{solute\_balance\_2d} does the same
in 2-D, uniform in $y$, and on a film open on all four sides with the flow
along the diagonal, which crosses an inlet--inlet, an outlet--outlet and two
inlet--outlet corners. The checks found three faults of the bookkeeping, not
of the fields; all corrected. \file{History.dat} counted only the solute the
flow carries out, not what the concentration equation diffuses through the
edge faces, and the balance stayed at $1.7\times10^{-4}$ of the outflowed
solute whatever $\Delta x$ and $\Delta t$ were. In 2-D it took the film and
the concentration in the boundary cell rather than at the face, a
first-order error in \code{VolErr} as well ($1.5\times10^{-3}$). With both
corrected it still counted what crossed from the velocities and the face
values of the field after each step, integrated by the trapezoidal rule,
not from the face fluxes the equations carried: \code{VolErr} and
\code{ResinVolErr} were $6\times10^{-7}$ here on 100 cells (and
$1.5\times10^{-6}$ with the open corners on 24 by 24), falling at order 1.8
in $\Delta x$. Counted from the step's own face fluxes they are
$1\times10^{-15}$ to $1\times10^{-14}$. \code{restart\_chain}
(\cref{sec:verification:bcio}) continues such a run from its state files, in
1-D and 2-D, to the bit of the run made straight through.

\code{nozzle\_balance} also runs \file{nozzle\_1d\_rotating} with its arm
turning at \SI{1}{\radian\per\second}: the nozzle crosses half the domain
and leaves it in part, and the film thickens from \SI{10}{\micro\metre} to
\SI{1.1}{\milli\metre} over cells \SI{3.3}{\micro\metre} wide while the
spinning substrate flings liquid into the $y$ that is not solved for. The
supply is about 200 times the initial volume, and \code{VolErr}, normalized
by the initial volume, must still be at rounding. The shipped case found it
at $-5.4\times10^{-5}$ ($2\times10^{-7}$ of the film; the turning arm
$4.3\times10^{-5}$), unchanged by tighter Newton and linear tolerances. The
supply and its bookkeeping agreed to rounding at every step; the volume the
step added did not, by $\Delta t$ times the sum over the cells of
$A\cdot\varphi$, which should be zero in a closed box. Formed as the product
of the entries of $A$ with $\varphi$, the capillary term's sum was
$\epsmach\,\lVert A\rVert\,\lvert\varphi\rvert$ per cell, entries of order
$D/\Delta x^{4}$ on a film 360 cells thick, and the edge faces leaked as
well: the 1-D \val{"Symmetry"} ghosts, two coupled rows of their own, came
off their mirror by rounding of the film. Corrected, the capillary and
vertical-gravity parts of $A\cdot\varphi$ are formed face by face from
differences of the film, the boundary rows are satisfied exactly after every
linear solve, and a \val{"Symmetry"} ghost is the exact mirror in 1-D as in
2-D (\cref{sec:time:apply}); \code{VolErr} is $2\times10^{-13}$ at the end
of the shipped case and $1\times10^{-13}$ at most with the arm turning. The
correction floor of the Newton iteration fell with it
(\cref{sec:convergence:floor}).

\code{nozzle\_supplied} repeats the three runs of \code{nozzle\_balance} with
the liquid supplied at a concentration of its own (\key{ConcType =
"supplied"}, \key{Conc} 0.9 into a film at 0.5, and 0.1 on the thick film):
\code{SuppliedResin} adds up $\Delta t\,S\,c_S$ from the very supply term the
solute balance took in, so \code{ResinVolErr} is at rounding
($5\times10^{-14}$), and the run fails the check if $c$ did not move. On the
thick film the supplied liquid at 0.1 lowers the viscosity under the nozzle
$e^{4}$-fold (\key{Exponent} 10), and in some steps the Newton correction
stagnates far from convergence, at $8\times10^{-4}$ to 1.9 times the first
correction. Such a step was found taken with a warning, the film keeping its
residual as a \code{VolErr} of $2.2\times10^{-8}$; corrected, it fails and is
cut (\cref{sec:convergence:stagnation}): 29 steps are cut, the run takes 3259
steps instead of 1982, and \code{VolErr} is $2\times10^{-13}$, checked on all
three films ($1\times10^{-14}$ on the closed boxes). The check fails if no
step of the thick film was cut, so that it keeps exercising the cut. \code{nozzle\_mixing}
supplies liquid at $c_S = 0.1$ onto a flat film at $c_0 = 0.5$ that does not
flow (no term on): with the nozzle wider than the domain, $h = h_0 + St$ and
$c$ follows \cref{eq:source-mixing}; under a nozzle of radius
\SI{0.1}{\milli\metre}, diffusion off, every cell keeps
$hc = h_0c_0 + (h - h_0)\,c_S$ with its own $h$. Each step adds $\Delta t\,S$
to $h$ and $\Delta t\,S\,c_S$ to $hc$ exactly ($S$ is constant), so there is
no error of the time discretization, and the largest error over 1-D and 2-D
is rounding, $4.2\times10^{-15}$. With \val{"local"} the same film keeps
$c = c_0$ exactly.

\code{nozzle\_arm} turns a nozzle on an arm over a film with no flow at all,
so that every point keeps what lands on it: the centroid of what arrived
between two outputs is where the nozzle was on average. The arm ramps up over
\SI{0.3}{\second}, turns for \SI{0.5}{\second} and ramps down over
\SI{0.3}{\second}. The ramp times were being multiplied by the time scale
$T_0$ instead of divided by it, so the arm moved on a clock of its own; the
check failed by \SI{30}{\percent} of the arm's length.

\code{nozzle\_area} holds a 2-D nozzle still in a closed box and compares
what it supplied with \key{InletRate} $\pi R^2 t$. The cells the nozzle
reaches were picked by the distance of their centers from it, up to
$R + 2\min(\Delta x, \Delta y)$, which on cells more than 3.9 times longer than
wide leaves out cells the circle covers in part: on the grid of cells 11 times
longer in $y$ the nozzle supplied \SI{81}{\percent} of its rate (found by
mutation testing, \cref{sec:notverified:mutation}). Corrected, the test allows a
whole diagonal of the cell, $R + \sqrt{\Delta x^2 + \Delta y^2}$, and what
is left is the error of counting $20\times20$ points a cell. On cells less
than 3.9 times longer than wide the results are the same, bit for bit.

\section{Manufactured solutions}
\label{sec:verification:mms}

The checks of \cref{sec:verification:analytic} use amplitudes of $10^{-3}$ of
the film, so the mobility $h^3$ hardly varies and the non-linear parts of the
discretization are barely exercised. The method of manufactured solutions has
no such limit. With \key{Numerics.ManufacturedSolution}
(\cref{sec:ref-numerics:mms}) the solver adds to the film-height equation the
source that makes
\[
  h = h_0 + a\cos(k_x x - \Omega t)\cos(k_y y)
\]
an exact solution, and \file{tools/mms.py} runs it with $a = 0.3\,h_0$ --- the
mobility varies by a factor of six over a wavelength --- with the capillary
pressure, vertical gravity and in-plane gravity all on, on a periodic domain
one wavelength across, for \SI{20}{\milli\second} at $\Omega =
\SI{60}{\per\second}$. The error is the largest difference from the
manufactured film at the end, relative to $a$; the order in $\Delta x$ is
measured at a time step too small to matter, the order in $\Delta t$ on a fine
grid against a run at a quarter of the smaller step. \code{ctest} runs it as
\code{mms}; it requires both orders to be at least 1.8.

The four rows of \cref{tab:mms} marked ``all'' switch on every term of the
model at once, each tuned to move $h_t$ (or $c_t$) by about as much as the others, so
that none hides behind another: over a substrate with a relief of
\SI{2}{\micro\metre}, the capillary pressure, both gravities, the
centrifugal term about an axis off the center of the domain, and --- in the
well-mixed model --- the disjoining pressure with the film at 1.75 to 3.25
times the precursor, the Marangoni flux, a viscosity, a surface tension and a
diffusivity that depend on the concentration, and the evaporation, with the
\val{"PowerN"} law in 1-D and Raoult's law with the \val{"negPow"} profile in
2-D. The concentration is manufactured as well,
$c = 0.5 + 0.1\cos(k_x x - 40t + 0.7)\cos(k_y y)$, and its equation gets its
own source, so its convection, diffusion and enrichment by evaporation are in
the comparison, and so is its error. With that many terms the derivatives
are too long to write by hand without a mistake --- and a mistake in the
source would look like one in the solver --- so \file{tools/mms\_generate.py}
derives them with SymPy from the equations as \cref{ch:governing} states
them, and writes them into \file{model/\bk ManufacturedSourceTerms.inc};
\code{ctest} checks, where SymPy is installed, that the file is what the
generator makes of the current equations (\code{mms\_source\_current}). The
evaporation, which needs no derivatives, is written out by hand in
\file{ManufacturedSource.cpp} from \cref{sec:terms:evaporation}, not taken
from the solver's own code.

\begin{yzntable}{The manufactured solution, error relative to $a$\label{tab:mms}}{@{}L{0.12\textwidth} L{0.40\textwidth} L{0.40\textwidth}@{}}
\textsf{\bfseries Case} & \textsf{\bfseries In $\Delta x$} & \textsf{\bfseries In $\Delta t$}\\ \midrule
\val{"WM"}, 1-D & $6.2\times10^{-3} \to 1.0\times10^{-3}$ (32, 64 cells): 2.61 & $9.5\times10^{-4} \to 2.2\times10^{-4}$ (1, \SI{0.5}{\milli\second}): 2.08 \\
\val{"WM"}, 2-D & $4.6\times10^{-2} \to 6.7\times10^{-3}$ (16, 32): 2.78 & $2.9\times10^{-3} \to 6.5\times10^{-4}$ (2, \SI{1}{\milli\second}): 2.15 \\
\val{"CZ"}, 1-D & $2.5\times10^{-2} \to 6.2\times10^{-3}$ (32, 64): 2.02 & $6.4\times10^{-3} \to 1.5\times10^{-3}$ (1, \SI{0.5}{\milli\second}): 2.06 \\
\val{"CZ"}, 2-D & $8.0\times10^{-2} \to 2.0\times10^{-2}$ (16, 32): 2.01 & $1.5\times10^{-2} \to 4.0\times10^{-3}$ (2, \SI{1}{\milli\second}): 1.94 \\ \midrule
\val{"WM"}, 1-D, all & $9.9\times10^{-4} \to 2.7\times10^{-4}$ (64, 128): 1.89 & $4.4\times10^{-4} \to 1.1\times10^{-4}$ (1, \SI{0.5}{\milli\second}): 2.07 \\
\val{"WM"}, 2-D, all$^{*}$ & $7.7\times10^{-3} \to 1.8\times10^{-3}$ (24, 48): 2.11 & $1.4\times10^{-3} \to 3.3\times10^{-4}$ (2, \SI{1}{\milli\second}): 2.07 \\
\val{"CZ"}, 1-D, all & $1.2\times10^{-2} \to 3.0\times10^{-3}$ (32, 64): 2.00 & $1.8\times10^{-3} \to 4.4\times10^{-4}$ (1, \SI{0.5}{\milli\second}): 2.05 \\
\val{"CZ"}, 2-D, all & $4.6\times10^{-2} \to 1.1\times10^{-2}$ (16, 32): 2.05 & $6.1\times10^{-3} \to 1.5\times10^{-3}$ (2, \SI{1}{\milli\second}): 2.01 \\ \midrule
\val{"CZ"}, 1-D, $c(x)$ & $7.4\times10^{-3} \to 1.9\times10^{-3}$ (32, 64): 2.00 & $4.0\times10^{-4} \to 8.5\times10^{-5}$ (1, \SI{0.5}{\milli\second}): 2.22 \\
\val{"CZ"}, 2-D, $c(x,y)$ & $2.9\times10^{-2} \to 7.4\times10^{-3}$ (16, 32): 1.97 & $1.6\times10^{-3} \to 3.2\times10^{-4}$ (2, \SI{1}{\milli\second}): 2.32 \\
\val{"CZ"}, 1-D, $c(x,\zeta)$ & $2.0\times10^{-2} \to 4.9\times10^{-3}$ (32, 64): 2.01 & $7.5\times10^{-4} \to 1.5\times10^{-4}$ (1, \SI{0.5}{\milli\second}; 128 cells): 2.30 \\
\val{"CZ"}, 2-D, $c(x,y,\zeta)$ & $7.7\times10^{-2} \to 2.0\times10^{-2}$ (16, 32): 1.96 & $6.4\times10^{-4} \to 1.6\times10^{-4}$ (1, \SI{0.5}{\milli\second}; 16 cells): 2.02 \\ \midrule
\val{"WM"}, 1-D, slip & $2.2\times10^{-3} \to 3.0\times10^{-4}$ (32, 64): 2.90 & $4.5\times10^{-4} \to 1.0\times10^{-4}$ (1, \SI{0.5}{\milli\second}): 2.12 \\
\val{"WM"}, 1-D, all, slip & $2.1\times10^{-3} \to 5.1\times10^{-4}$ (64, 128): 2.03 ($h$ 2.03, $c$ 1.90) & $6.6\times10^{-4} \to 1.6\times10^{-4}$ (1, \SI{0.5}{\milli\second}): 2.07 \\
\end{yzntable}
\noindent$^{*}$\,With a tenth of the diffusivity and half the relief of the 1-D case (see below).
The rows marked ``slip'' are \code{wm\_slip\_1d} and \code{wm\_all\_slip\_1d}: the
first two cases of their rows with Navier slip of half the mean film,
$\lambda = \SI{2.5}{\micro\metre}$ (\key{Physics.Slip}), which makes $\lambda h^2$
as large as $h^3/3$; the source is that of the slip mobilities
\cref{eq:q-slip}, generated by the same script.

The well-mixed orders above two in $\Delta x$ are the pre-asymptotic mixture
of a second-order capillary term and a third-order upwinded in-plane gravity,
whose error dominates on these grids; the \val{"CZ"} model takes both centered
and shows the second order plainly. Both models are second order in time with
every term on --- for \val{"CZ"} that is the predictor--corrector of
\cref{sec:time:newton}, without which the explicit terms of its film-height
equation would be first order. In the \val{"CZ"} rows marked ``all'' the
concentration is uniform. The rows marked $c(x)$ and $c(x,y)$ give \val{"CZ"} a
concentration that varies in the plane, $c = 0.5 + 0.2\cos(k_x x - 40t + 0.7)\cos(k_y y)$,
the same at every height in the film, with a viscosity that changes by
$e^{1.2}$ and a diffusivity by $e^{-0.8}$ over that range, $\sigma(c)$ and the
Marangoni flux, on top of every other \val{"CZ"} term. The concentration
equation gets its own source, $c_t + \mathbf{u}(\zeta)\cdot\nabla c -
\nabla\cdot(D\nabla c)$, which depends on $\zeta$ through the lubrication
profile; the error of $c$ is measured at every level, and the velocity
profile itself is compared with the exact one at every level. Each of $h$,
$c$ and $\mathbf{u}$ must converge at second order on its own (1-D: 2.00,
1.92, 2.01; 2-D: 1.97, 1.82, 2.00). The last two rows add a part that varies through the film,
$0.1\cos(k_x x-30t)\cos(k_y y)\cos\pi\zeta$, which meets the no-flux
conditions at the substrate and, without evaporation, at the surface. The
surface value then differs from the depth mean by 0.1 and the viscosity by
$e^{0.6}$ through the thickness. The diffusivity is
$1.5\times10^{-9}$~\si{\metre\squared\per\second}, a twentieth of that of
the rows above, so that the diffusion through the film stays within an order
of the other terms of $c_t$; it is not lower, because below about
$1.25\times10^{-9}$ in 2-D (and $10^{-9}$ in 1-D) the model does not keep
this state stable at short wavelengths (\cref{sec:verification:schemes}). Neither the flux nor the velocity profile
has a closed form: \file{mms\_generate.py} derives their integrands and the
mapped residual (the vertical advection $\Omega$ of the $\zeta$ grid, the
metric terms, diffusion at fixed $z$), and the source integrates them by
Gauss--Legendre quadrature; $N_z$ is refined with the grid. Each of $h$, $c$
and $\mathbf u$ must converge at second order in $\Delta x$ and in $\Delta t$
(1-D: 2.01, 1.98, 1.96 and 2.30, 2.04, 2.58; 2-D: 1.96, 1.88, 1.98 and
1.97, 2.02, 2.25). In 2-D the diffusion through the film,
$D\pi^2/h^2 = \SI{2.2e3}{\per\second}$ where the film is thinnest and $c$
lowest, makes a step of \SI{2}{\milli\second} 4.4 relaxation times long, and $\mathbf u$
converged at 1.7 on 2 to \SI{1}{\milli\second}; its order in $\Delta t$ is
therefore taken on 1 and \SI{0.5}{\milli\second}, on 16 cells, with 6 and 12
cells through the film on 16 and 32 cells (with 4 and 8, $c$ converged at
1.83 in $\Delta x$). The evaporation, whose surface condition $D c_z = cE$ the
manufactured profile does not meet, and the disjoining pressure (which
\val{"CZ"} does not have) stay out; the concentration through the
thickness is checked against a reference instead (\code{dry\_cz},
\cref{sec:verification:analytic}). The nozzle is outside the manufactured
solution in both models.

The runs with every term on found four faults that the runs with three had
no way to see:
\begin{itemize}
  \item The in-plane gravity and the centrifugal term took the viscosity at
    the concentration of the row's own cell, not of the face: the two rows
    that share a face saw different fluxes through it. The scheme was not
    conservative, and with a viscosity that depends on the concentration it
    did not converge at all (an error of 0.8 on every grid).
  \item Over a substrate with relief, the vertical gravity split each face's
    topography flux between the center and the neighbor and gave the
    center's share the viscosity of the cell rather than of the face ---
    the same fault, an error of $3.7\times10^{-2}$ that did not fall.
  \item The well-mixed concentration equation formed the explicit half of
    its first step before the first velocities had been computed, so the
    convection of that step was half at zero velocity: first order in
    $\Delta t$ overall ($1.24$), which the concentration then passed on to
    the film through the viscosity.
  \item In 1-D the well-mixed model counts the liquid the centrifugal term
    flings sideways, into the $y$ that is not solved for; \val{"CZ"} did
    not, so the two models disagreed on a spinning strip. Corrected,
    \val{"CZ"} counts it in the film-height equation and, through the continuity
    equation, in the vertical velocity.
\end{itemize}

The \val{"CZ"} runs with a concentration varying in the plane found five
more, all in the film-height equation and the velocity profile of
\val{"CZ"}:
\begin{itemize}
  \item The capillary term used the surface tension of the left $x$ face for
    the right face too, so again two rows saw different fluxes through the
    face they share (first order in $h$).
  \item The $x$ velocity of the centrifugal term ignored the $x$ offset of
    the rotation axis; the $y$ velocity and the film-height equation had it
    (a solute error of 0.19 on every grid).
  \item The first Simpson interval of the velocity integral took its
    midpoint at the substrate, which puts the profile off by
    $h^2\Delta\zeta^2/12\mu$ even at uniform viscosity (an error that did not
    fall).
  \item The 2-D capillary velocity lacked the cross derivatives
    $\partial_x\partial_{yy}h$ and $\partial_y\partial_{xx}h$ that the
    film-height equation has ($2.4\times10^{-2}$, not falling).
  \item At a periodic seam each face took the velocity of the face on the
    far side. For a periodic force the two are equal; the centrifugal force
    is not periodic, and the copy made the concentration first order.
    Corrected, each seam face keeps its own velocity, the one whose flux the film-height
    equation uses there.
\end{itemize}

The well-mixed model had the same seam fault: the far seam face took the
velocity of the near one, and with the centrifugal term the concentration
was first order (\code{wm\_all\_1d}: 1.24 on 32 to 64 cells, 0.92 on 64 to
128), which the combined order hid. Corrected, each seam face keeps its own
velocity; for a periodic force the two agree to rounding. \code{wm\_all\_1d}
requires $h$ and $c$ each at second order on its own (1.89 and 1.94, on
64 and 128 cells, with the default schemes). With the Kawamura--Kuwahara
scheme for the concentration (\cref{sec:verification:schemes}), the
concentration on 32 cells was still
pre-asymptotic, at 1.50: the third-order dissipation of its upwinded
convection, of the other sign from the rest of the error, took a part off
it; it reached 1.80 and 1.91 on the next two refinements (to 256 cells), and
1.96 to 1.99 with its convection switched off. In \code{wm\_all\_2d}, too, $h$ and $c$ must each
converge. The Kawamura--Kuwahara upwinding adds a third-order dissipation,
$-|u|\Delta x^3/4\,\partial^4$ (UTOPIA a third of it), of
the other sign from the second-order error of the diffusion (for $c$) and of
the capillary term (for $h$); near a cell Péclet number of $1/3$ the two
cancel, and the order stays below 2 on every grid a 2-D test can afford. With
the parameters of 1-D and KK for the concentration, $c$ converges at 1.66, 1.58 and 1.81 on 16, 32, 64 and
128 cells and $h$ at 1.96 and 1.89 on the last two pairs; with the upwind part
switched off (\val{"Central2"} for both), $c$ converges at 1.99, 1.98 and
1.95 and $h$ at 2.02, 2.00 and 1.96 on the same grids, and the error is smooth
everywhere, the seams and corners included. Switching the terms off one at a
time found no term that lowers the order by itself. The 2-D case therefore
uses a tenth of the diffusivity and half the relief, which puts the cell
Péclet number well away from the cancellation: on 24 and 48 cells $c$
converges at 2.58 and $h$ at 2.12 with KK for the concentration, and at
2.38 and 2.11 with the default schemes. The asymptotic order of the concentration itself,
past the cancellation, is checked in 2-D on a linear problem whose error is
known, an oblique sine carried and diffused at a lower cell Péclet number
(\code{oblique\_wm}, \cref{sec:verification:schemes}): second order with
UTOPIA and with Koren, settled on 16 to 64 cells.

The stratified \val{"CZ"} runs found four faults that a concentration uniform
through the thickness could not show:
\begin{itemize}
  \item The film-height flux took the surface tension and the viscosity at
    the depth-mean concentration, so it was not the depth integral of the
    velocity profile the concentration is carried by (errors of 0.2--0.3
    that did not fall). Corrected, it takes $\sigma$ at the surface
    concentration,
    and in the mobility and the Marangoni flux the means of $1/\mu(c(\zeta))$
    weighted by $3(1-\zeta)^2$ and $2(1-\zeta)$, which is that integral
    exactly (\cref{sec:disc:positivity-concz}).
  \item \code{Uz} was stored one cell too high.
  \item \code{Uz} was $-h\int\nabla\cdot\mathbf u\,d\zeta$ at fixed $\zeta$
    rather than the vertical velocity $w$, so the $\zeta$ velocity did not
    vanish at the surface of a sloping film and solute was carried through
    it.
  \item The ghost above the surface was not the image of the cell below it,
    which made the top cell's advection first order.
\end{itemize}
The field files also carried the velocity of the last step's midpoint, half
a step behind $h$ and $c$; corrected, they carry that of the state written.

\code{cz\_strat\_1d\_koren} runs the stratified case with the Koren scheme
for the concentration, which no other check does through the thickness. Its
orders in $\Delta t$ are taken on steps of 0.25 and \SI{0.125}{\milli\second}:
on the coarser pair, with the limiter settled at the midpoint of each
step, $\mathbf u$ is not yet in its asymptotic range (3.26; $h$ 2.22, $c$
2.05). \code{cz\_strat\_1d\_central4} runs it with Central4, which
dissipates nothing in any direction (\cref{sec:verification:schemes}).

With a concentration stratified through the thickness, the \val{"CZ"} step
was limited as $\Delta x^2$: on 128 cells a step of \SI{0.5}{\milli\second}
diverged. The cause was not the Marangoni term but two velocities that
disagreed. The film moved with the capillary and Marangoni flux at the
extrapolated midpoint concentration $c^*$, but the velocity the concentration
was carried with took its capillary part from the film halfway to the
\emph{predicted} end state, which had already balanced the Marangoni stress
of $c_n$. The difference, the Marangoni flux of $c^* - c_n$, left a
$\zeta$ velocity at the free surface that moved the stratification across its
own gradient: $c_t \approx D_M\nabla^2(c_n - c_{n-1})/2$ with
$D_M = h\sigma_c c_\zeta/2\mu$, stable only for
$\Delta t < \Delta x^2/2D_M$, which the thresholds observed on 64, 128 and 256
cells follow. Corrected, the velocity is formed from the film of the step
itself, and the coupling is iterated to the midpoint (\cref{sec:time:newton});
the step is limited only by accuracy. On 128 and 256 cells steps up to
\SI{2}{\milli\second} give ordinary time errors of a few $10^{-3}$, and an
evaporating 1-D film that had needed 1548 adaptive steps needs 134. The runs
in $\Delta t$ of \code{cz\_strat\_1d} are on 128 cells, and the
correction lowered the errors in $\Delta t$ of the stratified runs by a
factor of 4--5.

\section{Boundary conditions, restart and input parsing}
\label{sec:verification:bcio}

Three files cover this ground, and each fixes something different.

\file{test\_boundary\_condition.cpp} takes every combination of the five
boundary conditions --- \val{"Symmetry"}, \val{"Periodic"}, \val{"Pinning"},
\val{"FreshInlet"} and \val{"Outlet"} --- with the three fields
\code{Height}, \code{HeightJacobi} and \code{Conc}, and checks that the column
indices of the matrix stay within range, that every ghost row carries exactly
one equation (its diagonal is non-zero), and that a symmetric boundary
condition produces a matrix that is its own mirror image. It also checks
the values the conditions give: a uniform film is returned at every ghost
point, corners included, for every combination; an inlet on each edge holds
its face at the value of the field it was built from, whatever the interior
has become since; and mirroring or transposing the field together with the
edge types moves every ghost value with it. These found that an inlet on Xmax
or Ymax was held at the mean of its two ghost layers rather than of the two
points either side of its face, and that its face velocity was overwritten
with the one a cell inside, so that the solute carried in was not what the
film carried; both are fixed (the difference from the mirrored run fell from
$0.25$ to $2\times10^{-9}$). No shipped case has an inlet on those edges.

\file{test\_restart.cpp} exists because \file{field*****.dat} is both an output
file and the input format of an ASCII restart. \textbf{The line number written
and the grid point read back have to be the same point.} Three cases fix it:
(a) in 1-D, line $i$ lands at $x$ index $i+1$; (b) an asymmetric shape written
out and read back in returns to exactly the same places, to $10^{-13}$
relative; and (c) the decomposition of the line number is correct on a 2-D grid
as well. That last case uses a \textbf{$12 \times 5$} grid, and the asymmetry is
the whole point: \textbf{a mix-up between dividing the line number by $N_x$ and
dividing it by $N_y$ is invisible on a square grid and shows up only when the
two differ}.

\code{restart\_chain} (\file{tools/restart\_roundtrip.py --chain}) holds a
run that was stopped and restarted to one that ran straight through. Thirteen
cases of both models, in 1-D and 2-D, with the adaptive step, evaporation, the
coupling passes of \val{"CZ"}, a limited convection scheme, a moving nozzle
(once with its liquid at the local concentration, once at its own),
probes and the steady-state and dry-out stops, are stopped by \key{MaxSteps} at steps that
are not output times and restarted from the last file, alternating
\file{.vts} and \file{field*.dat}, into the same directory; one restarts from
an earlier file so that part of the run is done twice, two stop on the very
step at which the run becomes steady or dries out, and two (\code{wm1d\_pinned\_dry},
\code{\_adapt}) push a film against a \val{"Pinning"} wall and dry it, so that
the film at the wall passes a tenth of the thickest before the last
restart; two more (\code{wm1d\_open}, \code{wm2d\_open}) feed the film
through a \val{"FreshInlet"} and drain it through \val{"Outlet"} edges, so
that the totals of what crossed them go on from the state. Every field file, every row of
\file{History.dat}, \file{CalcInfo.dat} and the probes, and the step counts
summed over the parts must be identical, bit for bit. Without the state file
(\cref{sec:ref-numerics:restart}) two fixed-step cases must keep the output
times and agree to $10^{-11}$ (well-mixed) and $10^{-7}$ (\val{"CZ"}). The
check found five faults: after a restart the output schedule started again
from $t = 0$ and wrote a field every half interval; the totals of
\file{History.dat} started from zero; \val{"CZ"} took its first step from the
initial field, not the one it was given; the outer ghost column of a
\val{"CZ"} field read from \file{field*.dat} kept the initial film; and the
last step of each part was not counted toward the steady state, so a run
stopped on its steady step went on after a restart.

\file{test\_config.cpp} fixes the interpretation of the input files: scalars,
vectors, nested blocks and repeated blocks, and the requirement that a
\file{.toml} input produces the same tree as the equivalent
\file{.inp} (\cref{sec:inputformat:legacy}).

\section{Disjoining pressure and the contact line}
\label{sec:verification:disjoining}

This is the core of the chapter. The contact-line model is the part of Yuzen
with the most physics in it and the most ways to be subtly wrong, and it is
also the part where an analytic reference is available. The tests live in
\file{yuzen/tests/test\_disjoining.cpp}; the model they test is derived in
\cref{ch:wetting}.

\subsection{The calibration checked in closed form}

Equilibrium means uniform pressure, which for capillary plus disjoining
pressure is \cref{eq:eq-ode}, $H'' + \Dpn\,\Pihat(H) = K$ with $K$ the constant
pressure. In the limit where the precursor extends to infinity, $K = 0$, and
multiplying by $H'$ integrates the equation once:

\begin{equation}\label{eq:ver-first-integral}
  \tfrac{1}{2}\,(H')^{2} = -\Dpn \int_{\Hstar}^{H} \Pihat\,\mathrm{d}s .
\end{equation}

The integral to infinity is elementary (\cref{eq:pi-integral}), and it is
negative, so \cref{eq:ver-first-integral} gives the far-field slope
\cref{eq:slope-limit}: $(H')^{2} \to \As\,\thetae^{2}$. Multiplying the
dimensionless slope by $\varepsilon = h_0/L_0$ turns it into the physical
slope, and since $\As = 1/\varepsilon^{2}$ that physical slope is
\emph{exactly} $\thetae$. The coefficient the code builds,
\cref{eq:dp-def}, agrees with this identically; the derivation is in
\cref{sec:wetting:dp} and is not repeated here.

What the test adds is that the implementation really does evaluate that
formula, over a grid of parameters wide enough that an algebraic slip in any
one factor would show. It sweeps \textbf{36 combinations} ---
$\thetae \in \{2, 5, 10, 20\}^{\circ}$, $\Hstar \in \{0.01, 0.05, 0.2\}$ and
$(n,m) \in \{(3,2), (4,3), (9,3)\}$ --- and requires at every one of them that

\begin{itemize}
  \item $\Dpn$ agrees with \cref{eq:dp-def} to within $10^{-14}$ relative, and
  \item the recovered angle $\varepsilon\sqrt{-2\,\Dpn\int\Pihat}$ agrees with
    $\thetae$ in radians to within $10^{-12}$ relative.
\end{itemize}

A separate case fixes that disabling the disjoining pressure makes $\Dpn$
exactly zero. \textbf{This is the strongest claim in the chapter}: no
discretization and no time stepping are involved, and the tolerances are those
of double-precision arithmetic rather than of a numerical method.

\subsection{The residual velocity on the analytic equilibrium droplet}

Equilibrium droplets are a one-parameter family
(\cref{sec:wetting:equilibrium}): give $\Hmax$, solve \cref{eq:Hf-Hmax} for
$\Hf$ by bisection, and integrate $H'' = K - \Dpn\,\Pihat(H)$ from the foot to
the crest. No partial differential equation is involved.

The measurement is then simple to state. Put that shape on the mesh and add up
the capillary and disjoining-pressure velocities. In the continuous
equilibrium they cancel exactly, so \textbf{whatever is left is the truncation
error of the discretization}. The conditions are $\thetae = 10^\circ$,
$\Hstar = 0.02$ (that is \SI{0.2}{\micro\metre}), $\Hmax = 1.0$ (that is
\SI{10}{\micro\metre}), on a domain of $\pm\SI{300}{\micro\metre}$. Run the
solver, read \code{Ux} out of \file{field00000.dat} and divide by the velocity
scale $U_0 = 5.6069\times10^{-8}$~\si{\metre\per\second}.

\begin{yzntable}{Residual velocity on the analytic equilibrium droplet\label{tab:eqresidual}}{@{}d d d d d@{}}
\textsf{\bfseries $N_x$} & \textsf{\bfseries $\Delta x$ [\si{\micro\metre}]} &
\textsf{\bfseries residual velocity} & \textsf{\bfseries ratio} &
\textsf{\bfseries order} \\ \midrule
1000 & 0.6  & 18.94 &      &      \\
2000 & 0.3  & 5.074 & 3.73 & 1.90 \\
4000 & 0.15 & 1.294 & 3.92 & 1.97 \\
\end{yzntable}

The residual velocity in \cref{tab:eqresidual} is dimensionless. Taken from
$N_x = 1000$ to $N_x = 4000$ as a whole the order is \textbf{1.94}, and 1.97
on the finer pair: second-order convergence.

\textbf{The maximum sits at the contact line, where $H \approx 2\Hstar$ --- not
at the crest and not at the foot.} That is worth knowing for its own sake: it
says the truncation error of this model is concentrated exactly where the
profile bends most sharply, which is also where the mesh requirement of
\cref{sec:wetting:sizing} comes from.

The same quantity is measured a second time, independently, inside the test
itself. There the droplet is constructed on the spot with fourth-order
Runge--Kutta, the velocities of \code{LaplacePressure\_WM1D} and
\code{DisjoiningPressureKernel<WM1D\_Grid>} are added, and the maximum is
taken. Setting \env{YUZEN\_DBG\_EQ} to \val{1} prints the breakdown.

\begin{plain}
  [eq] nx=1000 umax=2.0295e+01  at X=-6.843746e-02 H=3.953799e-02
  [eq] nx=2000 umax=5.1789e+00  at X=6.869474e-02 H=3.599702e-02
  [eq] nx=4000 umax=1.3013e+00  at X=6.865186e-02 H=3.656075e-02
\end{plain}

The ratios there are 3.92 and 3.98. \textbf{The table produced by running the
solver and the independent integration inside the test agree to within
0.6 \% at $N_x = 4000$, 2 \% at $N_x = 2000$ and 7 \% at $N_x = 1000$} ---
two separate constructions of the droplet, two separate evaluations of the
velocity, and the residual is the small difference of two terms each about
170 times as large, so the two constructions' interpolation errors show in
it.

The assertion the test actually makes is looser than the numbers above: that
refining always reduces the residual, and that the ratio lies between 3.0 and
5.0. The band is wide because the interpolation error of the shape construction
rides along with the discretization error. It is nevertheless a sharp test,
because a wrong coefficient or a wrong sign puts the ratio near 1 rather than
near 4.

\subsection{Two traps in constructing the droplet}

\begin{warn}[Both of these destroy the convergence-order measurement itself]
  \file{test\_disjoining.cpp} carries both of these as comments in the code,
  and they are worth understanding rather than merely copying. Neither produces
  a slightly worse answer. \textbf{Each one makes the measured convergence
  order meaningless}, which is far harder to notice than a wrong number would
  be.

  \textbf{(a) The shape must be interpolated with cubic Hermite, not
  linearly.} Evaluating the second difference on the equilibrium shape divides
  by $\Delta x^{2}$. The error of a linear interpolation is
  $O(\Delta X^{2})$ in the step $\Delta X$ of the ordinary-differential-equation
  integration that built the shape --- which is \textbf{a fixed height error,
  independent of the mesh spacing $\Delta x$}. Divided by $\Delta x^{2}$ it
  therefore grows as $\Delta x^{-2}$, and \textbf{the residual velocity
  increases as you refine the mesh}. Cubic Hermite interpolation makes the
  error $O(\Delta X^{4})$, which leaves the refinement test alone.

  \textbf{(b) The initial value at the foot must be tuned so that the crest
  lands exactly on a grid point.} Scaling the foot's initial value $\eta_0$ by
  a factor $c$ shifts the whole shape along $X$ by $-\ln(c)/\lambda$, and that
  is the knob used to drive the last integration point onto the crest, where
  $H' = 0$. Skip it and the last point misses the crest by up to $\Delta X$, so
  the profile mirrored about the crest has a kink there. The height error of
  that kink is, once again, independent of the mesh spacing, so it destroys the
  convergence order for exactly the reason (a) does.
\end{warn}

\subsection{Time evolution from the equilibrium shape}

The equilibrium droplet at $N_x = 2000$ was then advanced in time for
\SI{0.01}{\second}. \Cref{tab:eqtime} gives what happens.

\begin{yzntable}{Evolution in time from the analytic equilibrium shape, $N_x = 2000$\label{tab:eqtime}}{@{}d d d d@{}}
\textsf{\bfseries $t$ [\si{\second}]} &
\textsf{\bfseries $\max\lvert\Delta h\rvert$ [\si{\micro\metre}]} &
\textsf{\bfseries relative} & \textsf{\bfseries residual velocity} \\ \midrule
0       & 0       & 0       & 5.074 \\
2.5e-3  & 4.23e-5 & 4.23e-6 & 0.709 \\
5.0e-3  & 5.36e-5 & 5.36e-6 & 0.418 \\
7.5e-3  & 6.02e-5 & 6.02e-6 & 0.296 \\
1.0e-2  & 6.48e-5 & 6.48e-6 & 0.229 \\
\end{yzntable}

The shape barely moves --- $6.5\times10^{-6}$ relative at the end --- and the
residual velocity \textbf{decreases monotonically}. That is the behavior to
expect, and the reason is worth stating: the shape it started from is the
equilibrium of the \emph{continuous} problem, and what it is relaxing toward
is the equilibrium of the \emph{discrete} one. The two are not the same shape,
and the difference between them is the truncation error measured in
\cref{tab:eqresidual}.

The apparent contact angle over the same interval goes from $9.56613^\circ$ to
$9.56610^\circ$: \textbf{unchanged to five digits}. The discrete equilibrium
it relaxes toward is the film with a uniform discrete potential, the minimum of
the discrete free energy (\cref{sec:disc:positivity-one}).

\begin{note}[Switch the steady-state test off to reproduce this table]
  With the default \key{SteadyStateCriterion} of \val{1e-3} the run is declared
  steady almost immediately and stops --- correctly, since the shape really is
  not changing. To produce \cref{tab:eqtime} you have to disable the test; see
  \cref{sec:verification:reproduce}, and \cref{sec:convergence:termination} for
  what the criterion means.
\end{note}

A related and simpler identity is checked in the same file. \textbf{A uniform
precursor film $H = \Hstar$ is an exact equilibrium}, because
$\Pihat(\Hstar) = 0$ identically. The test does not merely evaluate the
velocity there; it \textbf{dirties the velocity field first}, filling it with
the value 1.0, then calls \code{CalVelocity} and requires the disjoining
pressure's contribution to be below $10^{-12}$ at every interior \code{Xedge}.
Writing a zero is thereby distinguished from leaving a zero alone.

\subsection{The precursor bias}

A model with a \emph{finite} precursor does not deliver the angle you asked for
unless the droplet is thick compared with the precursor. That bias is derived,
tabulated and turned into sizing rules in
\cref{sec:wetting:apparent,sec:wetting:sizing}, and \cref{tab:apparent-angle}
is the table; none of it is repeated here.

What belongs in this chapter is the verification status of that table. It is
produced by \code{tools/make\_equilibrium\_1d.py --out /dev/null --sweep} at
$\thetae = 10^\circ$, $\Hstar = 0.02$, $(n,m) = (3,2)$, and two things about it
are measured rather than asserted. \textbf{The numerically integrated apparent
angle agrees with the closed form \cref{eq:theta-cap} to four digits.} And
\textbf{the maximum slope of the profile converges toward $\thetae$ far more
slowly than the apparent angle does} --- which is why the two columns of
\cref{tab:apparent-angle} have to be read separately, and why measuring an
angle off the steepest point of a computed profile needs a thicker droplet
again than measuring it by parabolic extrapolation.

\section{Spreading and Cox--Voinov}
\label{sec:verification:spreading}

\subsection{What is measured, and why it is not Tanner's law}

Tanner's law --- $R \sim t^{1/7}$ for a two-dimensional ridge --- is a law of
\emph{complete} wetting, $\thetae = 0$. This disjoining-pressure form cannot
represent that case at all, because by \cref{eq:dp-def}
\textbf{$\thetae \to 0$ implies $\Dpn \to 0$}: the term that produces the
contact line switches itself off in the limit. Tanner's law is therefore not
the right target here.

What can be measured in partial wetting is the relation that sits inside
Tanner's law, the \textbf{Cox--Voinov} relation

\begin{equation}\label{eq:cox-voinov}
  \theta_d^{3} - \thetae^{3} = 9\,\ln(L/l)\;\Ca,
  \qquad \Ca = \frac{\mu U}{\sigma}.
\end{equation}

Measuring it is cheap, and that is worth spelling out. \textbf{Each output time
of a single run is one $(\theta_d, U)$ point}, so a single calculation yields
the whole curve; there is no need to follow three decades of time to see the
relation.

The case is \file{cases/eq/spread\_n500.toml}: an equilibrium droplet built for
$\thetae = 20^\circ$ is used as the initial condition and the run is made with
\key{ContactAngle} set to \val{10.0}, so the droplet spreads toward its new
equilibrium shape. \file{tools/measure\_spreading.py} fits a parabola by least
squares to the outside of the droplet at each output time and extrapolates it
to $h = 0$, giving the base radius $R$ and the apparent angle $\theta$.

\begin{plain}
  theta : 19.1201 -> 9.4397 deg     (t = 0 .. 1 s)
  R     : 59.744  -> 85.507 um
  volume proxy theta*R^2 : 1.1911e-09 -> 1.2046e-09   (conserved to within 1.1 %)

  theta^3 - theta_e^3 = 19.6272 * Ca + 5.97e-05
    => 9 ln(L/l) = 19.627  =>  ln(L/l) = 2.181   (L/l = 8.9)
  correlation = 0.99791   (108 points, Ca = 1.41e-05 .. 1.12e-03)
\end{plain}

\textbf{The correlation is 0.99791 over 108 points spanning two decades of
$\Ca$.} The relation \cref{eq:cox-voinov} therefore holds in form, over a range
of capillary numbers wide enough that the fit is not an artifact of a short
lever arm. The intercept $5.97\times10^{-5}$ is about 1 \% of the scale of
$\theta^{3}$ --- $9.44^\circ$ is \SI{0.165}{\radian}, giving
$4.5\times10^{-3}$ --- which confirms that the equilibrium angle has been taken
correctly: a wrong $\thetae$ would show up as an intercept of the same size as
the signal.

\begin{warn}[The logarithmic factor is not a verified quantity]
  The fit gives $\ln(L/l) = 2.18$. The naive estimate --- the logarithm of the
  droplet half-width \SI{85.51}{\micro\metre} over the decay length of the foot
  \SI{1.196}{\micro\metre} --- is 4.27. These are the same order of magnitude,
  and that is the entire content of the comparison. \textbf{This logarithm
  shifts by a factor of about two depending on how the inner and outer lengths
  are chosen}, so the agreement claims nothing beyond order of magnitude. The
  value of $\ln(L/l)$ itself has not been verified.
\end{warn}

\subsection{The end point of spreading against an independent droplet}

The second half of the measurement is stronger, because it has a reference that
does not come from the solver. Let the spreading run all the way to
equilibrium, then compare the shape it reaches with an equilibrium droplet
\textbf{constructed without the solver}, from the ordinary differential
equation of \cref{sec:wetting:equilibrium}. Both are measured with the same
parabola fit, at \code{--hcut 1.0} \si{\micro\metre}.

\begin{yzntable}{End point of spreading against an independently constructed equilibrium droplet\label{tab:spreadend}}{@{}d d d d d@{}}
\textsf{\bfseries $N_x$} & \textsf{\bfseries $\Delta x$ [\si{\micro\metre}]} &
\textsf{\bfseries end of spreading [deg]} &
\textsf{\bfseries analytic droplet [deg]} & \textsf{\bfseries difference} \\ \midrule
500  & 1.2 & 9.4073 & 9.4052 & 0.02 \% \\
2000 & 0.3 & 9.4101 & 9.4083 & 0.02 \% \\
\end{yzntable}

The analytic droplet in \cref{tab:spreadend} is built to the same crest height
as the run it is compared with. The agreement is not only in the fitted angle
but in the shape itself: \textbf{the maximum difference between the two
profiles is \SI{0.0022}{\micro\metre} against a crest height of
\SI{7.0567}{\micro\metre} at $N_x = 2000$, that is 0.03 \%}
(\SI{0.0032}{\micro\metre} at $N_x = 500$).

The end point is the steady state of the scheme, the film whose discrete
potential is uniform (\cref{sec:disc:positivity-one}); outside the contact
line its shape is that of the continuous equilibrium even at $N_x = 500$,
where $\Delta x = \SI{1.2}{\micro\metre}$ is comparable with the transition
length $\ell = \SI{1.15}{\micro\metre}$ rather than a fifth of it. The path
there agrees as well: the apparent angle of the two meshes differs by at most
\SI{0.06}{\percent} over the first second. (With the arithmetic mean that the
disjoining pressure once took for its face mobility, the end point at
$N_x = 500$ was \SI{1.7}{\percent} off and that at $N_x = 2000$
\SI{0.11}{\percent}.) The rule $\Delta x \le \ell/5$ of
\cref{sec:wetting:sizing} stays a starting point: this droplet is a lenient
case, and refining is what tells.

\section{The cost of the coupling}
\label{sec:verification:coupling}

Thickness and concentration are solved as a block Gauss--Seidel iteration
inside the outer Newton loop (\cref{sec:time:coupling}). The question measured
here is how much of that outer iteration the concentration is responsible for.

The setup is \file{cases/eq/spread\_n500.toml} --- $N_x = 500$,
$\hstar = \SI{0.2}{\micro\metre}$, $\thetae = 10^\circ$, the same case verified
in \cref{sec:verification:spreading} --- with the concentration coupling turned
fully on: convection, diffusion, evaporation and Marangoni, with a viscosity
exponent of 10. \textbf{All conditions are compared at
\key{EndTime} $=$ \val{0.1}~\si{\second}}. Setting \env{YUZEN\_DEBUG\_RESIDUAL}
to \val{1} prints the thickness correction and the concentration correction for
every outer iteration, and the measurement counts the iterations on which the
concentration correction was the larger of the two.

\begin{yzntable}{Outer-iteration cost on the spreading family\label{tab:couplingspread}}{@{}L{0.25\textwidth} d d d d d@{}}
\textsf{\bfseries Condition} & \textsf{\bfseries steps} &
\textsf{\bfseries mean NK} & \textsf{\bfseries max NK} &
\textsf{\bfseries conc.\ dominated} & \textsf{\bfseries at final iter.} \\ \midrule
Evap.\ 1e-7, no disjoining pressure   & 1129 & 2.56 & 3 & 57.7 \% & 91.5 \% \\
Evap.\ 1e-7, with disjoining pressure & 1129 & 2.56 & 3 & 60.7 \% & 88.0 \% \\
Evap.\ 1e-6, with disjoining pressure & 1129 & 2.69 & 3 & 96.3 \% & 92.3 \% \\
Evap.\ 1e-5, with disjoining pressure & 1129 & 3.04 & 4 & 98.7 \% & 96.4 \% \\
\end{yzntable}

The evaporation coefficient of \val{1e-5} moves the concentration from a
maximum of 0.500 to 0.619 (mean 0.500 to 0.571) within the
\SI{0.1}{\second}, so as a coupling it is a genuinely demanding condition. The
stagnation cut-off of \cref{sec:convergence:tests} warned \textbf{zero times} in
all four.

Laying the corrections out iteration by iteration --- here the 1000th step of
the two runs at an evaporation coefficient of \val{1e-7} --- shows what is
happening.

\begin{plain}
  with disjoining pressure     thickness     2.62e-06 -> 7.23e-11 -> 1.84e-15
                               concentration 3.21e-06 -> 9.18e-10 -> 3.51e-14

  without disjoining pressure  thickness     6.20e-06 -> 1.38e-10 -> 3.47e-15
                               concentration 3.70e-06 -> 1.43e-09 -> 7.94e-14
\end{plain}

With the disjoining pressure on or off, the thickness falls by four to five
decades an iteration and is at rounding after two. The term is stiff --- $\Pihat \sim H^{-n}$
is the stiffest thing in the model (\cref{sec:terms:disjoining}) --- but its
Jacobian is exact (\cref{sec:verification:jacobian}), so it costs no extra
iterations; before that Jacobian was corrected the same runs took four to
five. Once the thickness has converged, the concentration correction is on
many iterations relatively the larger, and the more so the stronger the
evaporation: on \SI{61}{\percent} of the iterations at \val{1e-7} and on
\SI{99}{\percent} at \val{1e-5}. That, and nothing more, is what the
percentages in the table are.

\begin{warn}[Read the percentage the other way round]
  \textbf{``The fraction of iterations dominated by concentration'' measures how
  easy the height equation is, not how hard the concentration equation is.} The
  row with the highest figure in \cref{tab:couplingspread}, 98.7 \%, converges
  in 3.04 iterations on average, against 2.56 for the rows with the lowest. A
  large value is not a warning about the concentration solve.
\end{warn}

The same measurement was made on the cases that strengthen the coupling
deliberately, \file{cases/stress\_*.toml}, all of them derived from
\file{allphys\_1d\_reg}. Here too the stagnation cut-off warned zero times in all
four.

\begin{yzntable}{Outer-iteration cost on the \file{stress\_*} cases\label{tab:couplingstress}}{@{}L{0.19\textwidth} L{0.29\textwidth} d d d d@{}}
\textsf{\bfseries Case} & \textsf{\bfseries What is changed} &
\textsf{\bfseries steps} & \textsf{\bfseries mean NK} &
\textsf{\bfseries max NK} & \textsf{\bfseries conc.\ dom.} \\ \midrule
\file{stress\_evap}     & Evaporation coefficient $\times 100$ & 675 & 3.01 & 4 & 68.0 \% \\
\file{stress\_coupling} & Viscosity exponent 10 to 30; concentration coefficient of surface tension $\times 10$ & 675 & 2.19 & 4 & 2.2 \% \\
\file{stress\_dt}       & $\Delta t \times 100$, fixed (step control off) & 21 & 4.29 & 6 & 14.4 \% \\
\file{stress\_all}      & All of the above & 21 & 5.05 & 6 & 0.0 \% \\
\end{yzntable}

\subsection{What this measurement does and does not support}

\textbf{What it supports.} Over the eight conditions measured, the outer
iteration converged in \textbf{at most 6 iterations}, and the stagnation
cut-off never once fired. The fraction of iterations on which the concentration
correction governed the convergence test ranged over \textbf{0 to 99 \%}; at
the strongest evaporation, where it was highest, the mean was 3.04
iterations. Because the convergence test looks at both corrections
(\cref{sec:time:coupling}), the state it converges to is \textbf{the fully
implicit coupled solution of thickness and concentration together}. Writing the
off-diagonal blocks of the Jacobian would not change that answer; it would
change only the number of outer iterations, and that number is currently
between 2 and 6.

\textbf{What it does not support.} Three limits, each of which matters.

\begin{itemize}
  \item \textbf{This is a measurement of iteration counts, not of wall-clock
    time.} How much the cost of a single iteration rises when the coupling is
    switched on has not been measured at all.
  \item \textbf{Everything here is 1-D and \val{"WM"} (well-mixed).} It says
    nothing about two dimensions and nothing about the \val{"CZ"} model
    (\cref{ch:concentration}).
  \item \textbf{The step counts are those the adaptive time step chose}, not
    a comparison at a common step size. On the spreading family the step
    grew on its schedule without a cut under all four conditions, hence the
    same 1129 steps; the only thing held fixed by the set-up is the end time
    of \SI{0.1}{\second}.
\end{itemize}

\subsection{The coupling passes of \texorpdfstring{\val{"CZ"}}{CZ}}
\label{sec:verification:czpasses}

With the Marangoni term on, \val{"CZ"} solves the film and the concentration
of each step again, one stage each, at the midpoint concentration $c_m$ the
last pass found, until that midpoint settles
(\cref{sec:time:newton}). Each pass is a map $c_m \mapsto c_{\mathrm{mid}}(c_m)$,
and Aitken's relaxation takes $c_m + \omega\,(c_{\mathrm{mid}} - c_m)$ with
$\omega$ from the change of the residual between passes.

\textbf{The gain of a plain pass, derived.} Linearize about a flat film of
thickness $h$ at rest, with a profile $c_0(\zeta)$ through it, and perturb it
by a wave $e^{ikx}$: $\eta$ in the film, $\gamma(\zeta)$ in the
concentration. The film sees $c_m$ only through the Marangoni stress of its
surface value $\gamma_{m,s}$, and the concentration sees the film only through
the velocity in $\zeta$ that the Marangoni flow and the film's response to it
make, which carries the stratification $\partial c_0/\partial\zeta$. A plain
pass ($\omega = 1$) is therefore of rank one for each wave,
$\gamma_m \mapsto b + v\,\gamma_{m,s}$, and its one non-zero eigenvalue
$\lambda_k = v_s$ decides whether the passes converge. For a linear profile,
$\partial c_0/\partial\zeta = \Delta c$,

\begin{equation}\label{eq:cz-gain}
  \lambda_k = \frac{\Delta t}{2}\,D_M k^2\,\psi(1), \qquad
  D_M = \frac{h\,\sigma_c\,\Delta c}{2\mu},
\end{equation}

where $\psi$ solves
$\bigl(1 + \tfrac12\Delta t\,D k^2\bigr)\psi - \tfrac{\Delta t\,D}{2h^2}\,\psi''
 = \zeta(1-\zeta)\,\varphi(\zeta)$ with $\psi'(0) = \psi'(1) = 0$,
$\varphi = m_k - \frac{a}{1+a}\bigl(1 - \frac{\zeta}{2}\bigr)$ and
$a = \Delta t\,\sigma h^3 k^4/(6\mu)$, the Crank--Nicolson factor of the
capillary pressure. $\sigma_c$ is \key{Fluid.SurfaceTension.Coef}, $D$ the
diffusivity, $k^2$ the film's two-point symbol
$(2 - 2\cos k\Delta x)/\Delta x^2$, and $m_k$ the velocity's fourth-order
surface gradient over the film's two-point one (1 on long waves, $7/6$ on
the shortest). The first term of $\varphi$ is the Marangoni flow, the
second the return flow of the film it raises. Three properties follow.

\begin{itemize}
  \item $\lambda_k$ has the sign of $\sigma_c\,\partial c_0/\partial\zeta$
    and grows with $\Delta t$ and $k$. When the diffusion crosses the film
    within the step ($\Delta t\,D/h^2 \gg 1$) it tends to
    $\lambda_k \to \tfrac12 \mathrm{Ma}_h\,\bar s\,\theta/(1+\theta)$, with
    $\theta = \tfrac12\Delta t\,D k^2$, $\bar s = \int_0^1\zeta(1-\zeta)\varphi\,d\zeta$
    between $1/24$ (capillary pressure stiff) and $1/6$, and the Marangoni
    number of the film $\mathrm{Ma}_h = \sigma_c\,\Delta c\,h/(\mu D)$: the
    gain never passes $|\mathrm{Ma}_h|/12$ ($\times 7/6$ on the shortest
    waves; at most $0.072\,|\mathrm{Ma}_h|$ over a sweep of $\Delta t$ from
    $10^{-4}$ to $10^{2}$~\si{\second}, three film thicknesses, three diffusivities
    and 4 to 32 cells through the film).
  \item The fixed point of the passes is the implicit midpoint step of the
    linearized system $\dot y = A_k y$, and
    $1 - \lambda_k = \det(I - \tfrac12\Delta t A_k)\,/\,[(1+a)\det B_k]$, with
    $B_k$ the Crank--Nicolson factor of the diffusion through the film; both
    factors of the denominator are positive. So $\lambda_k > 1$ only where
    a wave of the model grows faster than $2/\Delta t$, a step on which the
    implicit midpoint rule turns that growth into a perturbation that
    changes sign. Where no wave grows, $\lambda_k < 1$ for every wave at
    every step.
  \item Aitken's factor settles to $\omega \approx 1/(1-\lambda)$ of the wave
    that holds the residual: between 0 and 1 for $\lambda < 0$, above 1
    for $0 < \lambda < 1$, negative for $\lambda > 1$. One wave, or a band
    of waves of about one gain, is then removed whatever its gain. Waves
    with gains on both sides of 1 cannot be removed by one factor together.
\end{itemize}

\file{tools/cz\_coupling.py} evaluates $\lambda_k$ of
\cref{eq:cz-gain} for every wave of a grid, for a linear profile, the
boundary layer of an evaporating film, or a column read from a field file,
together with the largest growth rate of $A_k$ and $1-\lambda_k$ from the
determinants. \env{YUZEN\_DEBUG\_COUPLING} prints every pass of a run.

\textbf{The gain measured.} An evaporating flat film (\SI{10}{\micro\metre}
at the start, \key{Physics.Evaporation.Coef} $= 3\times10^{-5}$ with a rate
that does not depend on $c$, $D = 10^{-9}$, 64 cells over
\SI{0.5}{\milli\metre}, \key{Grid.Nz} $=11$, a wave of 16 cells and
amplitude $10^{-3}$ of the film) ran \SI{0.04}{\second}, until
$\Delta c = 0.039$ to $0.042$ and $h = \SI{8.8}{\micro\metre}$; its last step
then took ten plain passes. After the first, which removes the drift of the
profile (gain 0), the residual is held by the shortest wave, and the ratio
of one pass's movement to the last approaches the largest $|\lambda_k|$ that
\file{tools/cz\_coupling.py} gives for the column of the run
(\cref{tab:czgain}): within 0.97 to 1.09 of it for gains of 0.34 to 2.4. At
\SI{20}{\milli\second}, with gains of 2 to 3, the ratio was still rising
after six passes. The signs are those of $\sigma_c$: on the same film at
\SI{10}{\milli\second}, Aitken's factor read $\omega = -0.64$ and $0.28$
for $\sigma_c = \pm0.05$, which is $\lambda = 1 - 1/\omega = 2.56$ and
$-2.57$, against $\pm 2.37$ predicted, and in nine passes the relaxed
ones took the residual down by $10^{4}$ and $10^{6}$ where the plain ones grew
it by 2.6 a pass.

\begin{yzntable}{Gain of a plain coupling pass on an evaporating flat film:
  largest $|\lambda_k|$ of \cref{eq:cz-gain} / measured ratio of successive
  passes\label{tab:czgain}}{@{}L{0.20\textwidth} L{0.17\textwidth} L{0.17\textwidth} L{0.17\textwidth} L{0.17\textwidth}@{}}
$\sigma_c$ [\si{\newton\per\metre}] & $\Delta t = \SI{2}{\milli\second}$ & \SI{5}{\milli\second} & \SI{10}{\milli\second} & \SI{20}{\milli\second} \\ \midrule
$+0.03$ & 0.336 / 0.358 & 0.836 / 0.843 & 1.419 / 1.529 & 1.955 / 2.636 \\
$-0.03$ & 0.336 / --- & 0.836 / --- & 1.419 / 1.539 & 1.955 / 2.602 \\
$+0.05$ & 0.560 / 0.542 & 1.393 / 1.392 & 2.365 / 2.589 & 3.250 / 4.322 \\
\end{yzntable}

(---: the residual reached rounding, about $10^{-11}$, before the shortest wave held it.)

\textbf{The passes the runs take.} On the stratified manufactured state of
\code{cz\_strat\_1d} (\cref{sec:verification:mms}), at a fixed step for
\SI{0.02}{\second}, the largest number of passes a step took was the same on
32, 64 and 128 cells: 1 at 0.5 and \SI{1}{\milli\second}, 2 at
2 and \SI{5}{\milli\second}, 4 or 5 at \SI{10}{\milli\second}, and at
\SI{20}{\milli\second} the passes did not settle (the midpoint still moved by
0.56 to 0.68 of the first pass's change after eight). The same holds with
$\sigma_c$ of the opposite sign and with a fifth of the diffusivity. What
fails at \SI{20}{\milli\second} is not \cref{eq:cz-gain}: with the viscosity
independent of $c$ (\key{Fluid.Viscosity.Exponent} $=0$ instead of 3) every
step on 32 cells settled in 2 passes. A viscosity that depends on $c$ couples the
concentration of every level to the velocity under the strong flows of that
state, and the pass is then no longer of rank one; a step of
\SI{20}{\milli\second} moves its wave by 1.2 radians. With four times the
$\sigma_c$, a film that is unstable in the model, the passes took up to 8 and
did not settle from \SI{10}{\milli\second}. An evaporating flat film under
step control (128 cells over \SI{0.5}{\milli\metre}, $\sigma_c = \pm0.01$ to
$\pm0.05$, \key{Criterion} $=10^{-2}$ and $10^{-1}$) took 1 to 3 passes a step,
and up to 8 where $\sigma_c > 0$; in two of those runs one step of 55 did not
settle. \file{concz\_1d\_reg} with evaporation up to
$3\times10^{-5}$, $\sigma_c$ up to $\pm0.03$ and 120 to 480 cells never took
more than 2.

\textbf{What happens when they do not settle.} Found: a step whose passes
did not settle within eight was kept as the last pass left it, with a
warning; the time-step control, whose truncation-error test looks at the film
before the passes, did not see it. On \code{cz\_strat\_1d} at 32 cells with
the step control on, \key{InitialDt} $=\SI{0.02}{\second}$ and a
\key{Criterion} of 1, loose enough to accept that step, the error at
\SI{0.02}{\second} was 2.1 of the film's amplitude; at a \key{Criterion} of
$10^{-3}$ to 0.3 the truncation-error test cut the step first. Corrected:
passes that do not settle fail the step, as a linear solve that does not
converge does. Under step control the step is cut by \key{ChangeFactor} and
solved again (the first such step is reported in \file{Warnings.txt}, and
\file{JobInfo.txt} counts them under \code{Steps failed: the passes did not
settle}); the run above then ends with an error of 0.17, that of the step
the truncation-error test chooses at a \key{Criterion} of 0.3. At a fixed
step the run stops, naming the passes as the cause. The midpoint also counts
as settled once it moves by less than ten times
\key{Numerics.ConcEquation.LinearSolver.Epsilon}, and no less than
$10^{-10}$, so that a step whose extrapolation was already exact is not
failed for the rounding of its solve. \code{coupling\_cz} of
\file{tools/features.py} checks both behaviors on that state.

\textbf{The safe regime.} Where the film is stable in the model, the plain
gain stays below 1 on every wave at any step; where $\sigma_c\,\partial
c/\partial\zeta < 0$ the gain is negative and the relaxation settles it at any
step, in a number of passes that grows with $|\lambda|$; and below
$|\mathrm{Ma}_h| \approx 10$ even the plain pass converges. In the runs
measured, the passes stopped settling only on films unstable in the model
($\sigma_c\,\partial c/\partial\zeta > 0$ with $\mathrm{Ma}_h$ of about
100 to 180, and the manufactured state with four times its $\sigma_c$)
and on steps on which a concentration-dependent viscosity moves the state by
about a radian; those steps now fail. The gain is derived for the
linearized flat film only; the viscosity's part is measured, not derived.

\section{The Jacobians}
\label{sec:verification:jacobian}

The Newton iteration of the film-height equation solves with the Jacobian
$J = \partial(A(h)h)/\partial h$, assembled as $A$ plus what each term adds for
the dependence of its coefficients on $h$. A mistake in that second part does
not change what the iteration converges to, only how fast; every other test
converges anyway, a little slower, and the first sign in a real run is a time
step that collapses. \file{test\_jacobian.cpp} therefore compares, term by
term, in 1-D and 2-D, $J v$ with the central difference $[R(h+\epsilon v) -
R(h-\epsilon v)]/2\epsilon$ of $R(h) = A(h)h$, on a film that is not flat, over
a substrate that is not flat, with a concentration that varies, for a
direction $v$ that is not smooth. What is left when the two agree is
rounding, $10^{-10}$; the test allows $10^{-7}$. It found three faults:

\begin{itemize}
  \item The vertical gravity and the disjoining pressure took the derivative
    of the face mobility $\tfrac12(h_o^3 + h_a^3)/\mu$ at the neighbor for the
    column of the cell itself as well ($2\times10^{-2}$).
  \item The disjoining pressure added all of $\hat\Pi'$ to a Jacobian that
    already held $\hat\Pi/h$ from $A$ (the form in which $\hat\Pi$ is written
    into the matrix): off by \SI{50}{\percent}.
  \item The sideways outflow of the 1-D centrifugal term added $2\,\mathrm{Sp}\,h/\mu$
    where $2\,\mathrm{Sp}\,h^2/\mu$ was due ($3\times10^{-2}$).
\end{itemize}

The evaporation is not in $J$ at all: the film-height equation takes it at
the current iterate, and its derivative with respect to $h$, which is zero
except near the precursor and at the cap, is left to the iteration.

The same file checks the face fluxes the well-mixed concentration is carried
with (\code{FaceFlux}, \cref{sec:disc:wm-conservative}): for every term and
every face interpolation of the upwinded ones, in 1-D (with the sideways
outflow of the centrifugal term) and 2-D, the signed sum of a row's face
fluxes must be the product of its entries of $A$ with $h$ to $10^{-13}$ of the
row, and the two cells that share a face must see the same flux to
$10^{-13}$ of the largest.

The three upwinded terms (in-plane gravity, centrifugal, Marangoni) are
checked with seven of the convection schemes of \cref{sec:disc:convection}.
The limited ones are checked on a film with steep fronts and no extremum,
which reaches every branch of Koren's limiter and stays at least 0.05 from
any of its kinks, with a finite-difference step of $10^{-7}$: van Leer's
curvature makes the difference itself $4\times10^{-7}$ off at $10^{-6}$.

\section{Choosing the convection scheme}
\label{sec:verification:schemes}

\file{test\_finite\_differences.cpp} fixes, for every scheme, the exactness
of the face value on polynomials (cubic for KK, UTOPIA and Koren where it is
not limited; quadratic for Central2, Central4, QUICK and Upwind2; linear for
van Leer), its order on a smooth profile, the bounds of the limited schemes on
random and step data and the overshoot of the linear ones, the bitwise mirror
and sign symmetry, the signs of the frozen weights and the derivative against
a finite difference. \file{tools/advection.py} (\code{ctest}: \code{advection})
carries a step in $c$ and a sine across a periodic flat film moved by in-plane
gravity, at a Courant number of 0.4, in both models; \file{tools/mms.py} and
\file{tools/equivalence.py} take \code{--conc}, \code{--height} and
\code{--scheme} to run any of their cases with another scheme. The defaults
were chosen from the comparison below, made with the other key at
\val{"KK"}.

\begin{yzntable}{Convection schemes on a concentration step and a sine (\file{tools/advection.py})\label{tab:schemes-advection}}{@{}L{0.16\textwidth} L{0.34\textwidth} L{0.40\textwidth}@{}}
\textsf{\bfseries Scheme} & \textsf{\bfseries Overshoot of the step, WM / CZ} & \textsf{\bfseries Sine, maximum error at 128 cells (order)}\\ \midrule
KK & \SI{6.3}{\percent} / \SI{7.0}{\percent} & $1.9\times10^{-4}$ (2.99) \\
UTOPIA & \SI{8.6}{\percent} / \SI{10.4}{\percent} & $6.3\times10^{-5}$ (2.99) \\
Central4 & \SI{36}{\percent} / \SI{39}{\percent} & $6.4\times10^{-4}$ (2.00) \\
Koren & 0 & $5.1\times10^{-3}$ (1.56; mean error 2.35) \\
van Leer & 0 & $1.0\times10^{-2}$ (1.41; mean error 1.98) \\
\end{yzntable}

The bound holds for a Courant number up to one only: Crank--Nicolson loses it
beyond, where Koren overshoots \SI{16}{\percent} at 1.5 and even first-order
upwinding \SI{9}{\percent} at 3. The solver reports a step past one
(\cref{sec:ref-numerics:convection}), and \file{tools/advection.py} checks
the report: with Koren, the front of \val{"WM"} is reported for the
concentration at 1.5 and not at 0.4, that of \val{"CZ"}, whose surface moves
at $1.5u$, at 1.5 and not at 0.4, and the film height, whose upwinded flux
moves at $3u$, at 0.4 and not at 0.3. (The film being flat, its report at
0.4 in the runs of \cref{tab:schemes-advection} has no consequence.)

The sine above is carried along $x$ without diffusion, so it measures the
convection alone. \file{tools/advection.py} also runs a check of the order
the well-mixed concentration reaches in 2-D with convection and diffusion
together (\code{oblique\_wm}, part of \code{advection}; it can be run alone
with \code{--oblique}). A flat square film, periodic in $x$ and $y$,
\SI{1}{\milli\metre} across, is carried by in-plane gravity at
\SI{1}{\milli\metre\per\second} in the direction $(2,1)/\sqrt5$, and its
concentration is a sine along the diagonal, $c = 0.5 + 0.2\sin k(x+y)$,
oblique to the flow and to both axes, so that neither the field nor its
motion is separable, with $D = 5\times10^{-7}$~\si{\metre\squared\per\second}.
The answer is $c = 0.5 + 0.2\,e^{-D|\mathbf k|^2t}\sin\mathbf
k\cdot(\mathbf x - \mathbf u t)$. The runs take 16, 32 and 64 cells a side
and 20, 40 and 80 steps over one decay time, \SI{25}{\milli\second}: the
time step is refined with the grid. The expected order follows from the
modified equation of the scheme. The central diffusion under-damps the wave
by $D(k_x^4+k_y^4)\Delta x^2/12$, UTOPIA over-damps it by
$(|u_x|k_x^4+|u_y|k_y^4)\Delta x^3/12$, and Crank--Nicolson adds an error in
$\Delta t^2$, here about \SI{1}{\percent} of the rest. So the error is of
second order, with a third-order part of the other sign whose share is
$r = (|u_x|+|u_y|)\Delta x/2D$, $0.67$ times the cell Péclet number. The
observed order is $2 + \log_2[(1-r)/(1-r/2)]$, below 2 until $r$ is
small. This is the cancellation that keeps the order of $c$ low in
\code{wm\_all\_2d} at the diffusivity of 1-D (\cref{sec:verification:mms}).
Here the cell Péclet number is 0.125 to 0.031 and the formula gives 1.94
and 1.97. Measured, UTOPIA converges at 1.91 and 1.97, with an error within
\SI{3}{\percent} of the formula's on every grid ($4.2\times10^{-3}$,
$1.1\times10^{-3}$, $2.9\times10^{-4}$ of the amplitude). Koren, which clips
the extrema of the sine to first order, converges at 1.88 and 1.92. The
check requires, for both, an order between 1.9 and 2.1 on 32 to 64 cells and
an order on 16 to 32 within 0.1 of it, that is, in the asymptotic range. The
convection is not lost behind the diffusion: first-order upwinding converges
at 0.72 and 0.87 here, and KK, with three times UTOPIA's dissipation, at
1.76 and 1.90. The film stays flat to rounding, so the check says nothing
about the film height, whose order in 2-D is that of the manufactured
solutions. The six runs take about \SI{20}{\second} on two cores.

\begin{yzntable}{Convection schemes in the manufactured solutions: error on the finer grid / order in $\Delta x$\label{tab:schemes-mms}}{@{}L{0.20\textwidth} L{0.13\textwidth} L{0.13\textwidth} L{0.13\textwidth} L{0.13\textwidth} L{0.13\textwidth}@{}}
\textsf{\bfseries Case} & \textsf{\bfseries KK} & \textsf{\bfseries UTOPIA} & \textsf{\bfseries Central4} & \textsf{\bfseries Koren} & \textsf{\bfseries van Leer}\\ \midrule
\multicolumn{6}{@{}l}{\emph{\key{Concentration} varied; the error of $c$}} \\
\code{wm\_all\_1d} & $8.3\times10^{-5}$ / 1.80 & $9.0\times10^{-5}$ / 1.94 & $9.1\times10^{-5}$ / 2.00 & $1.0\times10^{-4}$ / 2.21 & $1.3\times10^{-4}$ / 2.04 \\
\code{wm\_all\_2d} & $6.9\times10^{-4}$ / 2.58 & $4.8\times10^{-4}$ / 2.38 & $6.7\times10^{-4}$ / 1.80 & $2.7\times10^{-3}$ / 1.91 & $3.4\times10^{-3}$ / 2.21 \\
\code{cz\_conc\_1d} & $4.9\times10^{-4}$ / 1.72 & $5.5\times10^{-4}$ / 1.92 & $6.5\times10^{-4}$ / 2.00 & $5.2\times10^{-4}$ / 1.90 & $4.8\times10^{-4}$ / 1.94 \\
\code{cz\_strat\_1d} & $2.0\times10^{-3}$ / 2.01 & $1.9\times10^{-3}$ / 1.98 & $1.6\times10^{-3}$ / 1.98 & $2.0\times10^{-3}$ / 2.02 & $2.9\times10^{-3}$ / 2.19 \\
\code{cz\_strat\_2d} & $9.3\times10^{-3}$ / 1.78$^{\dagger}$ & $9.0\times10^{-3}$ / 1.88 & $7.2\times10^{-3}$ / 1.89 & $1.2\times10^{-2}$ / 2.04 & $1.8\times10^{-2}$ / 1.79 ($u$: 1.67) \\
\multicolumn{6}{@{}l}{\emph{\key{Height} varied; the error of $h$}} \\
\code{wm\_1d} & $8.1\times10^{-4}$ / 2.82 & $5.3\times10^{-4}$ / 2.26 & $1.0\times10^{-3}$ / 2.04 & $1.0\times10^{-3}$ / 2.61 & $3.0\times10^{-3}$ / 2.12 \\
\code{wm\_2d} & $5.5\times10^{-3}$ / 2.89 & $2.5\times10^{-3}$ / 2.50 & $3.4\times10^{-3}$ / 2.13 & $6.7\times10^{-3}$ / 2.78 & $1.3\times10^{-2}$ / 2.32 \\
\code{wm\_all\_1d} & $2.5\times10^{-4}$ / 1.85 & $2.7\times10^{-4}$ / 1.93 & $3.1\times10^{-4}$ / 2.00 & $2.7\times10^{-4}$ / 1.89 & $4.2\times10^{-4}$ / 2.01 \\
\end{yzntable}
\noindent$^{\dagger}$\,Pre-asymptotic, as KK is in \code{wm\_all\_2d} (\cref{sec:verification:mms}): on 24 and 48 cells $c$ converges at 1.93 ($h$ 2.02, $\mathbf u$ 1.96).

Every order in $\Delta t$ is between 1.9 and 2.3. The manufactured state of
\code{cz\_strat\_*} is one the model keeps stable. Where the surface tension
rises with $c$ and $c$ rises toward the surface ($\sigma_c\,\partial
c/\partial\zeta > 0$, over half of each wavelength here), a surface
concentration that varies along the film drives a Marangoni circulation whose
vertical velocity carries the stratification, and the viscosity's dependence
on $c$ feeds the change back to the surface; only the diffusion damps the
short waves this makes grow. A frozen-coefficient analysis of the linearized
equations (the lubrication profile with $\mu(c)$, the Marangoni stress, the
capillary pressure, the vertical advection of the stratification and the
diffusion, about the manufactured state at each point and time of the run,
for waves in every direction in 2-D; \file{tools/cz\_stability.py}) gives
the growth rate of each wavelength. At the diffusivity of the checks,
$1.5\times10^{-9}$~\si{\metre\squared\per\second}, every wave from a quarter
of the state's wavelength down to two cells of a grid twice as fine as the
finest of the checks (256 cells in 1-D, 64 in 2-D) decays, at
\SI{7.8}{\per\second} (1-D) and \SI{3.5}{\per\second} (2-D) or faster.

The instability was found with the stratified checks first made at
$3\times10^{-10}$~\si{\metre\squared\per\second}. There the analysis gives
growth rates of up to $7\times10^{2}$~\si{\per\second} for waves
\SI{2.6}{\micro\metre} to \SI{5}{\micro\metre} long --- one to three cells of
128 --- a factor of up to $e^{14}$ over the \SI{20}{\milli\second} of the
run; with $\sigma_c = 0$ or $\mu$ independent of $c$ the growth falls or
goes, and with the sign of $\sigma_c$ reversed it moves to where
$\partial c/\partial\zeta$ has the other sign. With Central4 or Central2,
\code{cz\_strat\_1d} converged in time to a solution with an oscillation two
to three cells long along $x$, strongest halfway through the film, the same
at steps of 0.125 and \SI{0.031}{\milli\second} and with 32 and 64 cells
through the film --- the growth of a solution, not of a scheme --- and $c$
converged at 0.48 in $\Delta t$; the dissipation of UTOPIA damped it on 128
cells and fewer, but on 256 cells the oscillation grew there too. (Central
schemes had stayed smooth earlier in development only while the layer
velocities of the concentration equation did not add up to the film's flux
on the shortest waves; once the volume fluxes were made to add up,
\cref{sec:disc:cz-conservative}, they grew.) Corrected by raising the
diffusivity: in 1-D every scheme converges at second order
(\cref{tab:schemes-mms}), Central2 as well ($c$ 1.90 in $\Delta x$, 2.05 in $\Delta t$), and on 128 and 256 cells
UTOPIA keeps its orders --- $h$, $c$ and $\mathbf u$ at 2.01, 2.05 and 1.99
in $\Delta x$, and at 2.31, 2.05 and 2.37 in $\Delta t$ on 256 cells ---
and so do KK and Central4, within 0.05 of these.
\code{cz\_strat\_1d\_central4} in \file{tools/mms.py} runs the 1-D case with
Central4 (\cref{sec:notverified:numerics}). A
limited concentration scheme clips the extrema that the no-flux conditions at
the substrate and the surface put into the profile through the film; in
\val{"CZ"}, where it is settled to the midpoint of the step, that costs up
to a factor of 2 in the error against UTOPIA, and in the well-mixed 2-D
case 6--7: it is kept for problems with real fronts. For the film
height the limiter costs little on a smooth solution (Koren's error is
1.2--1.3 times KK's) and the linear wave of \cref{sec:verification:analytic}
shows the same: its speed error at 32 cells per wavelength is
$6.7\times10^{-4}$ against $5\times10^{-5}$, falling at 2.85.

What it gains shows on a steep film over a thin precursor: a block
\SI{5}{\micro\metre} high sliding under strong gravity over a precursor of
\SI{0.1}{\micro\metre}, 100 cells, adaptive step. The lowest film height
reached is \SI{0.085}{\micro\metre} with KK, \SI{0.028}{\micro\metre} with
UTOPIA and \SI{0.0998}{\micro\metre} with Koren and van Leer, and Central4
stops with a step that underflows. Over a precursor of \SI{0.02}{\micro\metre}
KK, UTOPIA and Central4 all stop; Koren and van Leer finish with the precursor
intact, and at 200 cells KK still dips to \SI{0.007}{\micro\metre} where Koren
stays at \SI{0.01998}{\micro\metre}. The face height of a limited scheme lies
between its neighbors, which is what the positivity-preserving mobility of
\cref{sec:disc:positivity} assumes of it.

Every scheme passes the 33 checks of \cref{sec:verification:equivalence}
(worst $6.4\times10^{-9}$), runs every shipped case with the same number of
steps and cuts of the step as KK, and gives the same answer with 1 and 2
threads to the last bit. The Newton iterations are the same too, except in
\file{stress\_all} and \file{stress\_dt}, whose steps are a hundred times
longer: there Koren and van Leer for the concentration take 142 and 139
iterations against KK's 104 in \file{stress\_all}, and Koren 87 against 85 in
\file{stress\_dt}. The shipped cases differ
from KK by at most $1\times10^{-5}$ of $h$ with another \key{Height} and
$2\times10^{-6}$ with another \key{Concentration} (\file{stress\_all}:
$3\times10^{-5}$ and $1\times10^{-4}$); their concentration stays within
$0.5000$--$0.5020$, so they cannot tell the concentration schemes apart. Hence the defaults: \val{"Koren"} for \key{Height}, \val{"UTOPIA"} for
\key{Concentration} in both models. With these, every manufactured solution
passes (\code{wm\_all\_1d}: $c$ 1.94; \code{wm\_all\_2d}: $c$ 2.38, $h$ 2.11).

Comparing the schemes also found a fault in the Marangoni term: it was
upwinded by the sign of $\Delta c$ rather than of the flux, which with a
negative \key{Fluid.SurfaceTension.Coef} took the downwind side. No
shipped case has one.

\section{Features no other check switches on}
\label{sec:verification:features}

\file{tools/features.py} (\code{ctest} entry \code{features}) runs what the
cases never use. \code{probe\_wm\_1d}, \code{probe\_wm\_2d} and
\code{probe\_cz\_1d} put a probe near a point and require the last line of
\file{Probe00.dat} to be what the field file of the same instant holds at the
nearest point. The check found that the probes landed on the first point
within one cell of the one asked for rather than the nearest --- in 2-D the probe of the test
sat half a cell off --- and \val{"CZ"} compared the meters of the input with
its dimensionless grid, so its probes landed nowhere near where they were
asked for. They also check the position the probe file names in its header
and, in \val{"CZ"}, the level through the film; the header of a 1-D
well-mixed run named a point one cell short of the one its values came
from. All three were corrected.
\code{coupling\_cz} drives the coupling passes of \val{"CZ"} past their limit of
eight with a step of \SI{0.02}{\second} on \code{cz\_strat\_1d}: at a fixed step the run must
stop and name the passes, and under a step control loose enough to accept
that step otherwise the step must be cut, reported once and counted, and the
error stay below 0.5 of the film's amplitude (0.17; 2.1 with the step kept;
\cref{sec:verification:czpasses}). \code{stagnation\_wm} starts
\file{stress\_wetting}, whose first step at \key{InitialDt} stagnates far from
convergence (\cref{sec:convergence:stagnation}): at a fixed step the run must
stop and name the stagnation --- taken with a warning, the film went negative
on the next step and the run ended as diverged --- and under step control the
step must be cut until it converges (thirteen cuts), reported once and counted,
the run end normally and \code{VolErr} stay at rounding
($2\times10^{-15}$). \code{evap\_precursor\_wm} and
\code{evap\_precursor\_cz} run \file{leveling\_1d\_reg} and
\file{concz\_1d\_reg}, whose film thins below its precursor at the foot of
the substrate step, and require \code{EvaporatedVol} never to fall. The
check found the precursor correction $f_h$ negative below $\Hstar$
(\cref{sec:terms:evaporation}), which condensed solvent onto the film: over
\SI{1000}{\second} of \file{allphys\_1d\_reg} as much as \SI{43}{\percent}
of the initial volume, the mean concentration falling from 0.5 to 0.15.
Corrected, $f_h$ is zero there. \code{steady\_wm} and
\code{steady\_cz} require a leveling film
to stop with \code{Steady state} long before its \key{EndTime}, flat;
\code{dry\_wm} and \code{dry\_cz} a drying film to stop with \code{Dry out}
in the step in which the mean concentration passes \key{EndConcValue};
\code{maxsteps} the run to stop after exactly \key{MaxSteps} steps and say so;
and \code{adaptive\_dt} a run whose time step starts at \SI{1}{\nano\second}
to let the step grow and still end within the error it was asked for of a run
at a fixed small step ($1\times10^{-4}$ of the change, with
\key{Criterion} $=10^{-4}$). \code{adaptive\_clock\_*} cut and lengthen the step in both models,
in 1-D and 2-D, and require the reported time to be the sum of the steps
solved to $10^{-12}$; \code{adaptive\_source} that a relaxing nozzle supplies,
at the reported times, the volume its time function gives.
\code{restart\_edited\_wm} and \code{restart\_edited\_cz} raise one column of
a field file by \SI{1}{\percent} after the run wrote it: the restart must take
the field as it stands and warn that it is not the state saved beside it,
while an unedited copy continues without that warning
(\cref{sec:ref-numerics:restart}). \code{restart\_pinned} restarts a film
pushed against a \val{"Pinning"} wall and dried until the film at the wall is
\SI{14}{\percent} of the thickest: the run must go on, with the line of the
check of the pinned edges as a warning (\cref{sec:ref-bcout:boundary}).

\section{Input fuzzing and random inputs}
\label{sec:verification:fuzzing}

The readers of the input file (TOML and \file{.inp}) and of the restart files
(\file{field*.dat}, \file{.vts}, \file{field*.state}) were fuzzed with
libFuzzer under AddressSanitizer and UndefinedBehaviorSanitizer, about
\SI{2.7}{\hour} of one core and 14.5 million inputs, from corpora of the
shipped cases, their \file{.inp} forms and restart files of short runs
(\file{tools/fuzz\_corpus.py}; configure with \code{-DYUZEN\_BUILD\_FUZZERS=ON},
Clang only). A refusal is an accepted outcome; a crash, a sanitizer report or
a hang is not. They found: toml\texttt{++} reaching an unreachable branch on
a bare key with a non-ASCII character; the \file{.vts} readers trusting a
block length the grid did not have (a heap overflow), the ASCII readers
taking a short row with the missing values uninitialized, and the state
reader throwing on a huge array length; conversions that took trailing text,
\val{nan}, fractional counts, lists for scalars and dates without a word
(\cref{sec:inputformat:types}); and values that ran to nonsense --- a
negative \key{MaxDtFactor}, which never ended, a negative viscosity, a film
with no liquid, which ran with exit code 0 and NaN in \file{History.dat}.
All were corrected and are refused, and the well-mixed model stops on a
field that has broken down, as \val{"CZ"} does (\cref{sec:output:calcinfo}).

\file{tools/proptest.py} draws random valid inputs from a seed --- the model,
the dimension, small grids, any subset of the terms with parameters over
orders of magnitude, the edges, the initial shape, the convection scheme,
fixed or adaptive steps --- runs each for a few dozen steps and checks that it
ends cleanly with finite fields, a positive film and, for a bounded scheme at
a Courant number up to one, a bounded concentration; that the volume
($10^{-9}$) and the solute ($10^{-11}$ well-mixed, with its advection on;
$10^{-13}$ \val{"CZ"}) are kept, on a closed domain and, for the
well-mixed model, on one with \val{"FreshInlet"} or \val{"Outlet"} edges,
what crossed them counted, and that a uniform
concentration without evaporation stays uniform ($10^{-10}$ well-mixed,
$10^{-12}$ \val{"CZ"}); and, for a
part of them, that the transpose and the mirror give the transposed and
mirrored result, that one and two threads agree, and that a restart chain
gives the straight run to the bit. About 990 configurations were run. They
found: a failed concentration solve that left NaN in the starting guess of
every retry of the step; the corner of two \val{"FreshInlet"} edges, which
moved at every Newton iteration until the step collapsed; an extra step after
a restart at the step at which the film dried out; the divergence test of
\val{"CZ"} skipping the last step; a \file{.pvd} that dropped the entry of the
field a restart started from when its six-digit time was rounded up; and the
column labels of the \val{"CZ"} \file{CalcInfo.dat}. They then found the
faults that follow, all corrected before release.

\begin{itemize}
  \item The in-plane body forces and the Marangoni stress of the well-mixed
    model put liquid through a \val{"Symmetry"} or \val{"Pinning"} wall.
    Their flux is a face value of the film times a drive, not a gradient, and
    the face value was taken from the ghosts, which at a symmetric wall
    mirror the film. The wall-face velocity was zero; $A$, the Jacobian and
    the velocity were not. $|\code{VolErr}|$ reached $2.5\times10^{-4}$ in
    \code{wall\_balance\_2d}; corrected, it is $10^{-14}$. With it went a 1-D volume
    that took the ghosts into the edge cells, so that a pinned film seemed to
    lose $5\times10^{-5}$ of its volume, and a \code{VolErr} that did not
    count the nozzle. Nor did \code{ResinVolErr} count the solute the
    nozzle's liquid brings ($0.057$ in seed 1122, $1.9$ at the end of
    \file{nozzle\_1d\_fixed}); it is counted in \code{SuppliedResin}
    (\cref{sec:output:history}), and \code{nozzle\_solute} holds it
    to $10^{-12}$.
  \item A drying \val{"CZ"} film that was not flat lost 3 to
    \SI{25}{\percent} of the solute its evaporated fraction should have
    concentrated, whatever $\Delta t$, $N_x$ and $N_z$: the conditions at the
    surface and the substrate set only the vertical part of the diffusive
    flux, and the part along the slope leaked out. The concentration
    equation was also not conservative, and a uniform concentration drifted
    by $10^{-6}$, differently in a mirrored run. Corrected, it is solved in
    conservative form for $hc$ (\cref{sec:disc:cz-conservative}); the solute
    of a closed film is kept to $10^{-14}$, a uniform concentration stays
    uniform exactly, and mirrored and transposed runs agree to $10^{-12}$,
    with the Marangoni term too.
  \item The in-plane body forces of \val{"CZ"} crossed a \val{"Symmetry"}
    wall the same way ($2\times10^{-5}$, seed 1015; $3.7\times10^{-3}$ in
    \code{wall\_balance\_cz}). Every flux, every layer velocity and the
    in-plane diffusion of a wall face are zero after the correction, and the
    volume and the
    solute of a closed box are kept to $10^{-15}$.
  \item A drying film passed $c = 1$, and the cap deleted the solute above it
    (well-mixed \code{ResinVolErr} $-0.21$, seed 1206; \val{"CZ"}
    $4\times10^{-4}$, seed 1118; \code{dried\_wm} lost all of it,
    \code{dried\_cz} stopped). Corrected, evaporation is limited to the solvent there
    is (\cref{sec:terms:evaporation}), the solute is kept to $10^{-13}$ and $c$
    stays at most 1.
  \item When the time-step control lengthened the step after an accepted
    one, it did so before the clock was advanced: the clock added the new
    step, which had never been solved, and ran ahead --- by
    \SI{0.37}{\percent} at the end of \file{leveling\_1d\_reg}, up to
    \SI{0.5}{\percent} in the contact-line cases. The nozzle and the
    manufactured sources were evaluated at the wrong times, the steady-state
    rate was divided by the wrong step, and in \val{"CZ"} the film update and
    the concentration of that step used the new step while the film had been
    solved with the old. Corrected, the longer step applies from the next
    step.
  \item In 1-D the centrifugal term flings the film into the $y$ that is not
    solved for (\cref{sec:terms:centrifugal}), and \code{VolErr} and
    \code{ResinVolErr} did not count it: a spinning strip between walls
    showed what it lost as an error ($-0.12$ of the volume in
    \code{spin\_balance}). Corrected, the film-height equation's own
    sideways outflow, over the step it solved, is counted in
    \code{OutflowedVol} and \code{OutflowedResin}, and \val{"CZ"} writes the
    same balance columns (\cref{sec:output:history}); both balance to
    $2\times10^{-15}$.
  \item A film that dried to $c = 1$ against a \val{"Pinning"} wall (seed
    1035), continued from its state file, was not the run straight through:
    $c$ differed by up to $3\times10^{-3}$ an output later. The field, and
    with it the state file, held $c$ capped at 1, while the concentration
    equation of the run straight through went on from its own uncapped
    solution. The correction of seed 1206 above (the cap on the equation's
    own state, the evaporation limited to the solvent there is) removed it.
  \item The check of the pinned edges, which stops an initial condition whose
    film does not come down to the wall, was made again on a restart, and
    refused a film whose edge had evolved under the run's own physics:
    pushed against the wall and dried, it passed a tenth of the thickest film
    (seed 1206). The check is of the initial condition alone; on a restart
    its line is a warning (\cref{sec:ref-bcout:boundary}).
  \item A \val{"CZ"} film that evaporated a sizeable part of itself in one
    step changed its solute by $10^{-10}$ (seeds 1054 and 2105, against the
    $10^{-11}$ the property then allowed). The step puts the linear solver's
    residual back into the solute as a uniform shift of the change
    (\cref{sec:disc:cz-conservative}); the shift moves the solvent the
    surface term takes out as well, and the size it was given left a quarter
    of the residual in the solute. Across a \val{"Periodic"} seam, where each
    side forms the flux from its own ghost, the ghosts of the change held
    their copies only to the solver's residual, which the shift did not
    see. Corrected, the solute of the 147 closed \val{"CZ"} films among
    500 seeds is kept to $2.4\times10^{-14}$, the limited schemes included,
    and the property asks for $10^{-13}$.
  \item A \val{"Periodic"} edge was accepted with the centrifugal term, whose
    arm jumps by the length of the domain across the seam: not a periodic
    problem, with liquid created or lost at the seam. It is refused, except
    under a manufactured solution (\code{periodic\_centrifugal} in
    \file{tools/input\_guards.py}).
  \item An unbounded convection scheme of the well-mixed concentration
    (KK, UTOPIA, Central4, QUICK; Koren and van Leer past a Courant number
    of one) overshoots past $c = 1$ at a sharp front, and the cap at 1 on the
    committed state cut the solute above it: \SI{0.72}{\percent} of the
    solute as a front from $c = 0.99$ to 0.3 crossed a periodic domain of
    100 cells once with UTOPIA (falling only as $\Delta x^{0.7}$),
    \SI{3.0}{\percent} with Central4, \SI{0.23}{\percent} for UTOPIA on a
    drying front. Corrected, the solute above 1 is moved to the neighboring
    cells with room (\cref{sec:disc:convection}) and kept to rounding. The
    cap of the working field, which the coefficients see, also froze the
    limiter of Koren and van Leer at values that were not those of the
    iterate, and the advection lost its flux form there: 1 to
    \SI{2}{\percent} of the solute made over one crossing at a Courant
    number of 1.5. Corrected, the limiter
    is frozen at the iterate as solved; the solute is kept to the tolerance
    of the Newton iteration ($5\times10^{-9}$), and to rounding once the
    step is completed with the limiter settled at its answer
    (\cref{sec:disc:wm-conservative}). \code{conc\_overshoot} and
    \code{conc\_overshoot\_koren} (\cref{sec:verification:analytic}) hold
    both, and no shipped case overshoots.
  \item Through a \val{"FreshInlet"} or an \val{"Outlet"} the balances were
    held only to the error of the scheme: \file{History.dat} counted what
    crossed the edge from the velocities and the face values of the field
    after each step, not from the face fluxes the step had carried
    (\code{ResinVolErr} up to $1.1\times10^{-2}$ over seeds 1000--1099, in
    seed 1033, a 2-D film fed through an inlet against a wall). Corrected, the
    step's own face fluxes are counted (\cref{sec:output:history}), and the
    volume and the solute of an open domain are held as those of a closed
    one ($4\times10^{-15}$ over the same seeds; \code{inlet\_outlet},
    \code{solute\_balance}).
\end{itemize}

A further 100 seeds, run under AddressSanitizer, found that with the
time-step control off (\key{Numerics.\allowbreak TimeStepControl.\allowbreak Enable} $=$
\val{false}) a well-mixed step whose Newton
iteration failed --- here a linear solve that never met its tolerance at a
step far too long for the film --- was taken as the failed iteration left
it, the film it started from, and the run ended normally; its mirror image
diverged instead (seed 3094). Corrected, such a step stops the run with the
reason, since a fixed step has no shorter one to be retried at
(\code{wm\_fixed\_unsolved} in \file{tools/input\_guards.py}).

The concentration solves were found to do the same, with the step control
on as well as off: a linear solve of the well-mixed or the \val{"CZ"}
concentration that did not meet its tolerance left a warning, and its answer
was taken. The \val{"CZ"} film-height solve did the same at a fixed step.
Corrected, a concentration solve that fails fails the step: under step
control it is cut and solved again, in \val{"CZ"} from the concentration it
started from, and at a fixed step the run stops, naming the solve.
\file{tools/input\_guards.py} gives each model's concentration solver one
iteration at a tolerance of $10^{-30}$ at a fixed step
(\code{wm\_conc\_fixed\_unsolved}, \code{cz\_conc\_fixed\_unsolved}: refused,
naming the concentration solve), and one or two iterations without the ILU
preconditioner under step control, at a tolerance the first step cannot meet
and its cuts can ($10^{-8}$ and $10^{-10}$), from a first step the case's own
solver takes without a cut (\code{wm\_conc\_cut}, \code{cz\_conc\_cut}:
ten and eight steps cut, and the run ends normally). No shipped case has a
concentration solve that fails; all 28 give the same field files with the
correction as without it, bit for bit.

The checks of \code{proptest} were tightened with these: on a closed domain
the solute of \val{"CZ"} must be kept to $10^{-13}$, measured from the field
files, and a uniform \val{"CZ"} concentration must stay uniform to
$10^{-12}$.
\code{ctest} runs 20 small seeds and the fifteen of a regression list (the seeds that found faults, and 1282, a periodic film that dries) as
\code{proptest} (1035 and 1206 with the restart chain), in a few seconds; \code{--seed N --count M} runs more, and
\code{--only SEED --keep} reproduces one.

\file{tools/input\_guards.py} also restarts each model from a field file
with the concentration of one column set to 1.25 or to $-0.25$ --- outside
$[0, 1]$ on one side only --- and the run must stop at its first step as
diverged (\code{cz\_conc\_above}, \code{cz\_conc\_below},
\code{wm\_conc\_below}; the well-mixed model moves what lies above 1 to the
cells around after every step, so only the lower side applies to it). Both
models fill a cell below 0 from the cells around it after the solve, and the
well-mixed step mixes the column back into range by itself, so the test of
divergence takes the lowest $c$ of the start of the step and of the solve,
before the fill, as well as that of the field (found: the restarted well-mixed
run went on; corrected). A well-mixed film read \SI{1}{\micro\metre} below
zero in one column must stop the same way, the film thickness having gone
negative (\code{wm\_film\_below}): the run of a fixed step far too long for
the film, which reached that test before, now stops earlier, its Newton
correction having stagnated far from convergence (\code{wm\_diverges},
\cref{sec:convergence:stagnation}). It checks the limit on the size
of the grid, $(N_x+4)(N_y+4)(N_z+4) \le 5\times10^{7}$, at either side of it,
with \code{--check-input}, which runs nothing. A string of a \file{.inp}
file without its closing quote took the rest of the file as its value, and
every key after it kept its default unless it was required; it is refused,
naming the key and the line (\code{inp\_unterminated}).

\section{Compilers and optimization}
\label{sec:verification:compilers}

The shipped cases (capped at 300 steps) and \code{ctest} were run with GCC 13
and Clang 18 at \code{-O0}, \code{-O2} and \code{-O3}, with and without
\code{-fno-strict-aliasing}, with uninitialized locals set to a pattern and
to zero, with scrambled heap memory (\env{MALLOC\_PERTURB\_}), and with the
checked standard library (\code{-D\_GLIBCXX\_DEBUG}). Every one gives the same
field files, to the bit, and GCC and Clang agree with each other. GCC 11,
the system compiler of RHEL 9, builds it as well and passes \code{ctest}
(it found a header that used \code{size\_t} without including
\code{<cstddef>}, which the newer compilers forgave; corrected). Fused
multiply-add (\code{-march=native}, or any arm64 build) and Eigen's vector
width change $h$ by about $10^{-12}$ of its maximum on ordinary cases. On the
contact-line cases the difference follows the floor at which the Newton
correction stagnates and can reach $10^{-8}$ over thousands of steps;
velocities at equilibrium are residuals and may differ by a few percent. The
solver's own code gives the same bits on any instruction set; only the
linear solver's does not. \code{-ffast-math} runs every case normally, but
three unit tests fail, and \code{-funsafe-math-optimizations} rewrites the
powers of the disjoining pressure less accurately, which shows at
$N_x = 4000$: do not build with it. The comparison found an exact-equality
unit test that failed under \code{-ffp-contract=fast} because two paths of
the face interpolation were fused differently (corrected: the derivative returns
the value's own result), and that \code{(int)} of a NaN read from the input is
undefined. A scan for what MSVC's OpenMP 2.0 rejects (unsigned loop indices,
\code{collapse}, array or min/max reductions, \code{simd}) and other
non-standard constructs found none; GCC with \code{-Wall -Wextra -Wpedantic
-Wvla} warns about nothing.

\section{Memory, threads and floating point}
\label{sec:verification:sanitizers}

None of the checks above would notice a read past the end of an array that
happens to find a plausible value, a race between threads that happens to
resolve the same way, or a NaN that the time-step control quietly retries
away. These are looked for separately; \cref{tab:sanitizers} gives what was run
and what it found before the release. CI builds with AddressSanitizer and
UndefinedBehaviorSanitizer on every push and runs the quicker solver tests
in that build, since a fault of this kind changes no answer that a test of
the answers could see (\cref{sec:notverified:mutation}).

\begin{yzntable}{The checks for memory, threads and floating point\label{tab:sanitizers}}{@{}L{0.24\textwidth} L{0.44\textwidth} L{0.24\textwidth}@{}}
\textsf{\bfseries Tool} & \textsf{\bfseries Run on} & \textsf{\bfseries Found}\\ \midrule
ASan and UBSan (GCC, assertions on) & the unit tests, every test tool including the whole of \code{mms} and 100 more seeds of \code{proptest}, and every case in \file{cases/} with each convection scheme, 20 steps & see below \\
Valgrind memcheck & the unit tests, every case in \file{cases/} and the scheme variants, three steps each, the \val{"CZ"} coupling passes, restarts with and without \file{field*.state}, and drying films at the solvent limit & no error, no leak \\
TSan (Clang, with the Archer tool of LLVM's OpenMP) & the unit tests, 1-D and 2-D runs of both models with every convection scheme, the \val{"CZ"} coupling passes, closed walls, evaporation to the solvent limit and restarts, two to four threads & no race \\
\code{feenableexcept} for invalid, divide-by-zero and overflow & every case in \file{cases/} with each scheme, and every run of the test tools & see below \\
clang-tidy (\code{bugprone}, \code{clang-analyzer}, \code{performance}), GCC with \code{-Wshadow -Wconversion -Wduplicated-cond -Wnull-dereference} & every source file & no defect \\
Coverage (gcovr) & the whole suite and every case & \SI{90}{\percent} of the lines; see below \\
\end{yzntable}

The checks were first run before the convection schemes and the coupling
iteration of \val{"CZ"} were added, and repeated after. The first time, the
floating-point traps fired in three places. The sigmoid that shapes the
initial film overflows \code{exp} far from its edge and returns the right
limit, 0 --- harmless. Eigen's BiCGSTAB breaks down on a few linear systems of
\file{stress\_wetting} and reports it, which the Newton iteration
handles by shrinking the step. And that is where the third was: the retried
step built the explicit half of its update from the attempt it had just
rejected, whose film had reached $10^{76}$, so it failed at once and the step
was cut again for nothing. That was a fault (\cref{sec:verification:equivalence},
\code{wm\_retry}), and it is fixed; every rejected step in every case had
been retried from the wrong film, which moved the results by about
$10^{-7}$.

The second time the traps fired in two places, both harmless: the sigmoid
again, and at a pinned wall where the film height is zero the capillary
velocity of the wall face, which is 0/0 and which the boundary condition
overwrites before anything reads it. Tested on its own, the van Leer limiter
was faulty: it formed $2ab/(a+b)$ and its derivatives from the products,
which underflowed to $0/0$ for differences below about $10^{-162}$ and
overflowed above $10^{154}$. Corrected, it uses the ratios $a/(a+b)$ and
$b/(a+b)$, which changed van Leer results only in the last bits. The coverage of the
code added since the first round (\SI{97}{\percent} of its lines) showed that
no test drove the \val{"CZ"} coupling passes to their limit of eight;
\code{coupling\_cz} in \file{tools/features.py} was added to do so. The field files are
bit-identical at one, two and four threads with every scheme and with the
coupling passes.

The third time followed the conservative form of the \val{"CZ"}
concentration, the wall faces of both models, the limit of evaporation to the
solvent, the clock and the restart state file. UBSan found that two unit
tests copied an \code{Input} whose flags had never been set, which is
undefined (in the tests only; corrected). The traps fired at the sigmoid of
the initial shapes and at the 0/0 velocity of a pinned wall face, in 2-D
as well, as in the earlier rounds; at the deliberate overflows of the input checks (a
viscosity exponent tested for overflow, a 3000-digit number in a restart
file); and in attempts of a step that the time-step control then rejected,
or in runs that stop as diverged --- never in an accepted step. The 100
further \code{proptest} seeds found the fixed-step fault of
\cref{sec:verification:fuzzing}. \SI{95}{\percent} of the lines added in
this round are covered; what is not is mostly the messages of damaged
restart files and refusals, and a field file edited after it was written,
for which \code{restart\_edited\_*} was added to \file{tools/features.py}.
The field files are bit-identical at one, two and four threads for the 2-D
well-mixed case with Koren and van Leer, and for \val{"CZ"} in 2-D with the
Marangoni coupling passes, closed walls and evaporation.

The fourth time followed the face-by-face form of $A\cdot\varphi$ and the
boundary rows satisfied exactly (\cref{sec:time:apply}). TSan found no race
at two and four threads (2-D well-mixed runs with Koren and van Leer, closed
boxes with \val{"Pinning"} and \val{"Symmetry"} edges and gravity, \val{"CZ"}
in 2-D), Valgrind and the sanitizers nothing, and 100 further \code{proptest}
seeds under AddressSanitizer no violation. The traps fired at the sigmoid of
the initial shapes and at the 0/0 velocity of a pinned wall face --- of the
vertical gravity as well as of the capillary pressure --- which the boundary
condition overwrites. A unit test of the boundary rows found that the two
rows of an \val{"Outlet"} or a \val{"FreshInlet"}, each of which refers to
the other's ghost, were left off by a ninth of what the solve had left, and
those of a 1-D pinned concentration 729 times further off; corrected, each
set of rows that refer to each other is solved together, which changes no
shipped case. \SI{99}{\percent} of the lines of this code are covered; the
field files are bit-identical at one, two and four threads for the 2-D
well-mixed case with Koren and van Leer, for \val{"CZ"} in 2-D and for a
pinned 2-D box.

The fifth time was the last before the release, on the input file in the form
this manual documents. AddressSanitizer and UndefinedBehaviorSanitizer ran
the unit tests, \code{input\_guards}, \code{error\_paths},
\code{restart\_roundtrip}, \code{features} and nine checks each of
\code{equivalence} and \code{analytic}; Valgrind 20 steps of
\file{tutorial\_1d} and \file{leveling\_2d\_reg}; TSan with Archer the unit
tests at two and four threads and 2-D runs of both models with either backend;
and the three fuzzers ran ten minutes each, 6.2 million inputs. AddressSanitizer
found that a run whose \file{History.dat} was a device that never ends
(\file{/dev/full}, in \code{error\_paths}) read it whole to resume it, which
took all the memory there was before the run could report that it cannot
write there; corrected, only a regular file is read. Nothing else was found.

To repeat them: configure a build with
\code{-DCMAKE\_CXX\_FLAGS="-fsanitize=address,undefined"} and
\code{-DYUZEN\_PREFER\_LIBOMP=OFF} (LLVM's \file{libomp} keeps its thread
structures to the end of the process, which LeakSanitizer reports as leaks)
and run \code{ctest} in it; run \code{valgrind --error-exitcode=77 build/bin/yuzen case.toml} with
\key{Numerics.MaxSteps} set small; build with \code{clang++
-fsanitize=thread} against \code{libomp} and run with
\env{OMP\_TOOL\_LIBRARIES} pointing at \file{libarcher.so}. Without Archer,
TSan reports the OpenMP runtime's own synchronization as races. For the
traps, a small \code{LD\_PRELOAD} library that calls \code{feenableexcept}
in every thread (through \code{pthread\_create}) and logs each trap site
before letting the run continue is enough.

\section{Long runs and large grids}
\label{sec:verification:long}

The checks above take seconds to minutes each. Two questions need longer runs.
Do the balances of \file{History.dat} stay at rounding over hundreds of
thousands of steps, or do they drift? And does memory grow with the step
count? \Cref{tab:longruns} answers both for the shipped cases run much
longer than they are, single-threaded. \code{VolErr} and \code{ResinVolErr}
are the largest over every row of the run. The resident memory is no larger
at the end than a quarter of the way through.

\begin{yzntable}{Long runs: the balances and the memory\label{tab:longruns}}{@{}L{0.30\textwidth} r r L{0.13\textwidth} L{0.13\textwidth} r@{}}
\textsf{\bfseries Run} & \textsf{\bfseries Steps} & \textsf{\bfseries $t$ [s]} & \textsf{\bfseries $|$VolErr$|$} & \textsf{\bfseries $|$ResinVolErr$|$} & \textsf{\bfseries Memory}\\ \midrule
\file{allphys\_1d\_reg}, $\Delta t = 10^{-4}$\,s fixed & 500{,}001 & 50   & $1.1\times10^{-14}$ & $2.3\times10^{-14}$ & 8.0\,MB \\
\file{allphys\_1d\_reg}, adaptive                       & 100{,}860 & 1000 & $8.6\times10^{-15}$ & $9.7\times10^{-15}$ & 7.9\,MB \\
\file{concz\_1d\_reg}, $\Delta t = 10^{-4}$\,s fixed    & 200{,}001 & 20   & $4.0\times10^{-15}$ & $5.3\times10^{-15}$ & 8.1\,MB \\
\file{allphys\_2d\_reg}, adaptive                       & 3{,}612   & 30   & $1.6\times10^{-14}$ & $1.8\times10^{-14}$ & 14.5\,MB \\
\end{yzntable}

In the adaptive 1000\,s run the film loses \SI{43}{\percent} of its volume,
\SI{26}{\percent} to evaporation and \SI{17}{\percent} through the sideways
outflow of 1-D (\cref{sec:output:history}), and its mean concentration rises
from 0.5 to 0.71; the balance still closes to rounding. No step failed in any of
the four runs. The fifth run writes a field file at every step, 30{,}001 of
them. The \file{.pvd} that lists them stays valid, and its time stamps stay
distinct; what it costs is given in \cref{sec:output:vtk}.

Three grids far larger than any in the test suite were each run for a few
steps, to measure the memory a step needs. The peak resident memory was:
\begin{itemize}
  \item well-mixed 2-D at $1024^{2}$ cells (\file{scaling/bench2d\_n1024}),
    \SI{3.8}{\giga\byte} with \val{"eigen"} and \SI{3.2}{\giga\byte} with
    \val{"amgcl"};
  \item well-mixed 1-D at $N_x = 10^{6}$, \SI{2.0}{\giga\byte};
  \item \val{"CZ"} 2-D at $200 \times 200 \times 20$ (\file{concz\_2d}),
    \SI{1.6}{\giga\byte}.
\end{itemize}
That is about \SI{3.7}{\kilo\byte} per cell in 2-D, \SI{2.0}{\kilo\byte} in
1-D and \SI{2.1}{\kilo\byte} per cell of $N_x N_y N_z$ in \val{"CZ"}. These
runs measure memory and nothing else: whether the linear solver converges as
well on such grids, and how the rounding accumulates over a long run on them,
is not measured (\cref{ch:notverified}).

\code{python3 tools/long\_runs.py --exe build/bin/yuzen --workdir runs/long}
repeats all eight runs in about an hour on one core. It prints the table,
with the resident memory a quarter of the way through and at the end.

\section{Reproducing every number in this chapter}
\label{sec:verification:reproduce}

Every command below is run from the root of the repository.

\subsection{The test suite}

This covers \cref{sec:verification:suite,sec:verification:grid,%
sec:verification:linalg,sec:verification:velocity,sec:verification:bcio} and
the calibration and equilibrium-droplet tests of
\cref{sec:verification:disjoining}.

\begin{shell}
cmake -S . -B build -DCMAKE_BUILD_TYPE=Release
cmake --build build -j
ctest --test-dir build                        # all seventeen entries
./build/bin/yuzen_tests                       # case and assertion counts
./build/bin/yuzen_tests -sf="*test_grid.cpp"  # the per-file breakdown
YUZEN_DBG_EQ=1 ./build/bin/yuzen_tests -sf="*test_disjoining.cpp"
\end{shell}

The last line prints the in-test residual velocities of
\cref{sec:verification:disjoining} --- 42.996, 10.801 and 2.7116. To run that
one case rather than the whole file, use
\code{-tc="*residual velocity on the equilibrium shape*"}.

\subsection{Thread count and backend}

For \cref{sec:verification:linalg,sec:verification:backends}:

\begin{shell}
./tools/run_case.sh cases/leveling_1d_reg.toml runs/reg_th1 1
./tools/run_case.sh cases/leveling_1d_reg.toml runs/reg_th2 2
python3 tools/compare.py runs/reg_th1 runs/reg_th2            # => 0.000e+00

./tools/run_case.sh cases/allphys_2d_reg.toml runs/2d_th1 1
./tools/run_case.sh cases/allphys_2d_reg.toml runs/2d_th2 2
python3 tools/compare.py runs/2d_th1 runs/2d_th2
\end{shell}

\file{tools/run\_case.sh} is used throughout the recipes below
(\cref{sec:env:tools}).

\Cref{tab:wholesolver} is \val{"eigen"}. On a build with AMGCL, \val{"auto"}
resolves to \val{"amgcl"} for the 2-D \file{allphys\_2d\_reg}, and the second
pair then may differ in the low-order bits
(\cref{sec:verification:amgclthreads}); write \val{"eigen"} into both
\key{Backend} keys as below, or configure with \code{-DYUZEN\_WITH\_AMGCL=OFF},
to reproduce the table.

To change the backend, write \val{"eigen"} or \val{"amgcl"}
into \key{Numerics.HeightEquation.\allowbreak LinearSolver.Backend} and
\key{Numerics.ConcEquation.\allowbreak LinearSolver.Backend} and run the same
commands. \textbf{Both keys have to be written out}: no case under
\file{cases/} sets \key{Backend} at all, so each equation would otherwise take
the \val{"auto"} default and resolve it for itself. The
total BiCGStab iteration count is the sum of the \code{BCG\_iter} column of
\file{CalcInfo.dat}.

\val{"amgcl"} needs a binary that has it. The build above leaves
\code{YUZEN\_WITH\_AMGCL} at its default of \val{ON}, but a build that could
not find or fetch the headers has silently turned it off
(\cref{sec:install:options}), and so has any build configured with
\code{-DYUZEN\_WITH\_AMGCL=OFF}; either way the run falls back to
\val{"eigen"} with a warning, and you measure \val{"eigen"} twice. So
configure it explicitly and check the \code{Solver Info} section of
\file{JobInfo.txt} afterwards:

\begin{shell}
cmake -S . -B build -DCMAKE_BUILD_TYPE=Release -DYUZEN_WITH_AMGCL=ON
cmake --build build -j
ctest --test-dir build                        # all seventeen on this build too
\end{shell}

\subsection{The equilibrium droplet}

The ASCII restart file is produced by \file{tools/make\_equilibrium\_1d.py}.
The files under \file{cases/eq/} agree with the output of this command byte for
byte. Adding \code{--sweep} prints the apparent-angle table,
\cref{tab:apparent-angle}, and writes no file --- but \code{--out} is a
required argument and is parsed before the sweep branch, so it still has to be
given, and \file{/dev/null} is the obvious thing to give it.

\begin{shell}
python3 tools/make_equilibrium_1d.py --out /dev/null --sweep

python3 tools/make_equilibrium_1d.py --out cases/eq/eq_n2000.dat \
        --nx 2000 --xmin=-0.3e-3 --xmax=0.3e-3 \
        --hstar 0.02 --theta 10 --hmax 1.0

for n in 1000 2000 4000; do
  ./tools/run_case.sh cases/eq/eq_n$n.toml runs/eq_n$n 1
done
\end{shell}

Take the largest absolute value of the third column (\code{Ux}, in
\si{\metre\per\second}) of \file{field00000.dat} in each output directory and
divide it by the velocity scale $U_0 = L_0/T_0$, with
$T_0 = \mu L_0^{4}/(\sigma h_0^{3})$, to get the dimensionless residual
velocity. For this case

\[
  L_0 = 1.749042\times10^{-3}~\si{\metre}, \quad
  T_0 = 3.119463\times10^{4}~\si{\second}, \quad
  U_0 = 5.606869\times10^{-8}~\si{\metre\per\second} .
\]

The scales are printed in the run log of every run as well
(\cref{sec:output:running}).

The time evolution uses \file{cases/eq/eq\_run.toml}, but the default
steady-state test fires immediately after the start, so add
\key{SteadyStateCriterion} $=$ \val{-1} to the \code{[Numerics]} table to
switch it off. Setting \key{EndTime} to \val{0.0102} gives the five \file{field*.dat} at
$t = 0, 2.5\times10^{-3}, \dots, 1.0\times10^{-2}$, and a sixth at the end time. The apparent angle is
measured with \code{fit\_cap} below.

\subsection{Spreading}

\begin{shell}
./tools/run_case.sh cases/eq/spread_n500.toml runs/spread 2
python3 tools/measure_spreading.py runs/spread/spread_n500 --dt-out 1e-3 --stride 5
\end{shell}

\code{--dt-out} must match \key{Output.FieldData.TimeInterval} in the input.
\code{--stride 5} uses every fifth file; the wider interval over which
$U = \mathrm{d}R/\mathrm{d}t$ is differenced reduces the noise in $\Ca$.

\begin{note}[The correlation depends on the stride]
  The numbers in \cref{sec:verification:spreading} are those of
  \code{--stride 5}. The correlation is 0.984 at \code{--stride 1}, 0.998 at
  \code{--stride 5} and 0.999 at \code{--stride 10}. It is a property of the
  differencing as well as of the physics, so quote the stride with it.
\end{note}

For the comparison of \cref{tab:spreadend}, extend \key{EndTime} to \val{8.0}
and let the run stop on its own steady-state test --- $t = \SI{2.23}{\second}$
at $N_x = 500$ and $t = \SI{2.24}{\second}$ at $N_x = 2000$. Read the crest
height from the last \file{field*.dat}, build an equilibrium droplet of the
same crest height with
\code{make\_equilibrium\_1d.py --theta 10 --hmax <crest height/h0>}, and put
both through \code{fit\_cap()} of \file{measure\_spreading.py}.

\begin{plain}
import sys; sys.path.insert(0, "tools")
from measure_spreading import fit_cap
R, theta, hmax = fit_cap("runs/spread/spread_n500/field02232.dat", 1.0)
\end{plain}

\subsection{The cost of the coupling}

Copy \file{cases/eq/spread\_n500.toml} and set the following.

\begin{itemize}
  \item \key{Physics.ConcField.EnableConvection} and \key{EnableDiffusion} to
    \val{true};
  \item \key{Physics.SolutoCapillarity.Enable} to \val{true};
  \item \key{Physics.Evaporation.Enable} to \val{true}, and
    \key{Physics.Evaporation.Coef} to \val{1e-7}, \val{1e-6} and \val{1e-5} in
    turn;
  \item \key{Physics.DisjoiningPressure.Enable} to \val{true} and to
    \val{false};
  \item \key{Numerics.EndTime} to \val{0.1}.
\end{itemize}

\begin{shell}
YUZEN_DEBUG_RESIDUAL=1 ./build/bin/yuzen case.toml 2> resid.txt
\end{shell}

Split the \code{[resid] <iteration> <thickness correction> <concentration
correction>} lines of \file{resid.txt} into time steps at each iteration number
0, and count on what fraction of lines the concentration correction exceeded
the thickness correction. The \file{stress\_*} cases of
\cref{tab:couplingstress} use their inputs exactly as they stand.
